\chapter{局部化}\label{chap:localization}\index{局部化|hyperpage}
\SourcePage{298}
本章定义等变上同调，并说明强有力的局部化方法。这里的思想与技巧，不仅将完成\hyperref[chap:gromov-witten]{第七章}开始的部分工作，还将提供\hyperref[chap:mirror-theorem]{第十一章}证明镜像定理所需的一些工具。

先在\hyperref[sec:localization-theorem]{第~\ref*{sec:localization-theorem} 节}介绍等变上同调的定义和基本性质，并讨论局部化定理 \SourceCite{AB} 及其各种推论。基本思想是，当光滑簇具有群作用时，簇上的上同调类“局部化”为固定点轨迹上的类，从而容易开展计算。将会看到，固定点集的法丛的等变 Euler 类起着重要作用。作为应用，先描述 $\mathbb P^r$ 的局部化，再用它确定一般五次三维簇上的直线数。就计数几何而言，局部化首次由文献 \SourceCite{ES2} 利用 Hilbert 概形上的环面作用引入。

本章主要工作从\hyperref[sec:stable-map-localization]{第~\ref*{sec:stable-map-localization} 节}开始，将局部化应用于模空间 $\overline M_{0,n}(\mathbb P^r,d)$。射影空间 $\mathbb P^r$ 有代数环面 $T=(\mathbb C^*)^{r+1}$ 的自然作用，从而诱导 $T$ 在该模空间上的作用。按照 \SourceCite{Kontsevich2}，将学习如何用某些带标记的图 $\Gamma$ 描述固定点集的分支，以及如何计算相应固定点分支的等变法丛。然后，在\hyperref[subsec:multiple-cover-localization]{第~\ref*{subsec:multiple-cover-localization}} 和\hyperref[subsec:nodal-quintic-double-covers]{第~\ref*{subsec:nodal-quintic-double-covers} 小节}继续\hyperref[subsec:instanton-numbers]{第~\ref*{subsec:instanton-numbers} 小节}开始的 Calabi--Yau 三维簇 Gromov--Witten 不变量研究。先计算三维簇中刚性嵌入的光滑有理曲线的 $d$ 重覆盖所作的贡献，再专门考察五次三维簇，研究五次结点有理曲线的二重覆盖这一更微妙的问题。这样将严格证明\hyperref[chap:gromov-witten]{第七章}关于瞬子数 $n_{10}$ 的论断。

本章最后一节简要介绍等变 Gromov--Witten 不变量与虚拟基本类的局部化。

\section{局部化定理}\label{sec:localization-theorem}
先定义\emph{等变上同调}。这里给出拓扑构造；代数构造参见 \SourceCite{EG}。设 $G$ 为连通 Lie 群，其分类空间由主 $G$ 丛 $EG\to BG$ 给出，总空间 $EG$ 可缩。此 $G$ 丛在同伦等价意义下唯一确定。我们将使用代数环面 $G=(\mathbb C^*)^n$；此时 $BG=(\mathbb{CP}^\infty)^n$，可取
\[
EG=\pi_1^*S^\times\times_{BG}\cdots\times_{BG}\pi_n^*S^\times,
\]
其中 $\pi_i:BG\to\mathbb{CP}^\infty$ 为第 $i$ 个投影，$S$ 是 $\mathbb{CP}^\infty$ 上截面层为 $\mathcal O_{\mathbb{CP}^\infty}(-1)$ 的重言丛，$S^\times$ 表示去掉零截面。
\SourcePage{299}
\subsection{等变上同调}\label{subsec:equivariant-cohomology}
若拓扑空间 $X$ 具有 $G$ 作用，令 $X_G=X\times_GEG$；它是 $BG$ 上纤维为 $X$ 的丛。为以后使用，令 $i_X:X\hookrightarrow X_G$ 表示一条纤维的包含映射。
\begin{definition}\label{def:equivariant-cohomology}
$X$ 的\emph{等变上同调}定义为
\[
H_G^*(X)=H^*(X_G),
\]
其中 $H^*(X_G)$ 为 $X_G$ 的通常上同调。
\end{definition}
\index{等变上同调|hyperpage}
等变上同调具有通常上同调的许多性质，例如存在平坦等变拉回与固有等变推前。这些运算满足通常上同调中相应运算的常见性质 \SourceCite{AB}。注意 $H_G^*(\mathrm{point})=H^*(BG)$。沿 $X\to\mathrm{point}$ 拉回可见，一般的 $H_G^*(X)$ 是 $H^*(BG)$ 模，因此可将 $H^*(BG)$ 看作等变上同调的系数环。此外，$i_X:X\to X_G$ 诱导“遗忘”映射 $i_X^*:H_G^*(X)\to H^*(X)$。

主要考虑 $G=(\mathbb C^*)^n$ 的情形。令 $M(G)$ 为环面 $G$ 的特征标群。每个 $\rho\in M(G)$ 给出一维向量空间 $\mathbb C_\rho$，$G$ 通过 $\rho$ 作用于其上。若 $L_\rho=(\mathbb C_\rho)_G$ 为 $BG$ 上的相应线丛，则 $\rho\mapsto c_1(L_\rho)$ 定义同构 $\psi:M(G)\simeq H^2(BG,\mathbb Z)$，进而诱导复系数环同构 $\operatorname{Sym}(M(G))\otimes\mathbb C\simeq H^*(BG,\mathbb C)$。称 $\psi(\rho)$ 为 $\rho$ 的\emph{权}，以下未注明系数的上同调均取复系数。

特别地，若 $\rho_i$ 是由 $\rho_i(t_1,\ldots,t_n)=t_i$ 定义的 $G=(\mathbb C^*)^n$ 的特征标，令 $\lambda_i$ 表示其权。于是得到同构
\begin{equation}\label{eq:equivariant-coefficients}
H_G^*(\mathrm{point})=H^*(BG)\simeq\mathbb C[\lambda_1,\ldots,\lambda_n].
\end{equation}
将线丛 $L_{\rho_i}$ 记为 $\mathcal O(\lambda_i)$，则 $\lambda_i=c_1(\mathcal O(\lambda_i))$。在式~\eqref{eq:equivariant-coefficients} 的等同下，$i_{\mathrm{point}}^*$ 可看作将全部 $\lambda_i$ 映到 $0$ 的“非等变极限”。更一般地，对任意 $X$，映射 $i_X^*$ 都可看作非等变极限，特别会将全部 $\lambda_i$ 映到 $0$。

一般而言，对于 $G$ 空间 $X$，\emph{等变向量丛}是 $X$ 上的向量丛 $E$，使 $G$ 在 $X$ 上的作用提升为 $E$ 上的作用，并在纤维上为线性作用。此时 $E_G$ 是 $X_G$ 上的向量丛；\emph{等变 Chern 类} $c_k^G(E)\in H_G^*(X)$ 定义为通常 Chern 类 $c_k(E_G)$。若 $E$ 的秩为 $r$，最高 Chern 类 $c_r^G(E)$ 称为 $E$ 的\emph{等变 Euler 类}，记为
\[
\operatorname{Euler}_G(E)\in H_G^*(X).
\]
\hyperref[chap:gromov-witten]{第七章}中的 Euler 类是同调类，即最高 Chern 类的 Poincaré 对偶；从这里开始，将通常或等变 Euler 类看作上同调类更为方便。
\index{等变向量丛|hyperpage}\index{Chern 类!等变|hyperpage}\index{Euler类@Euler 类!等变|hyperpage}
\SourceExampleNumber{9.1.1}{1}
\begin{example}\label{ex:diagonal-equivariant-chern}
$G=(\mathbb C^*)^n$ 在 $\mathbb C^n$ 上的对角作用，给出 $X=\mathrm{point}$ 上的等变向量丛 $E$，满足 $E_G\simeq\mathcal O(\lambda_1)\oplus\cdots\oplus\mathcal O(\lambda_n)$。因此，$\lambda_1,\ldots,\lambda_n$ 是该\IndexLocation{indexloc:1434:299}表示的权。因为 $\lambda_i=c_1(\mathcal O(\lambda_i))$，有
\begin{equation}\label{eq:equivariant-chern-symmetric}
c_k^G(E)=\sigma_k(\lambda_1,\ldots,\lambda_n)\in\mathbb C[\lambda_1,\ldots,\lambda_n],
\end{equation}
其中 $\sigma_k$ 为第 $k$ 个初等对称函数。
\end{example}
\SourcePage{300}
\subsection{Atiyah 与 Bott 的定理}\label{subsec:atiyah-bott-localization}
接下来需要局部化定理。设环面 $T=(\mathbb C^*)^n$ 作用于紧复流形 $X$。由 \SourceCite{Iversen}，固定点轨迹 $X^T$ 是光滑连通分支 $Z_j$ 的并。令 $i_j:Z_j\hookrightarrow X$ 为包含映射，$N_j$ 为 $Z_j$ 在 $X$ 中的法丛。$N_j$ 是等变向量丛，故具有等变 Euler 类
\[
\operatorname{Euler}_T(N_j)\in H_T^*(Z_j).
\]
等变包含映射 $i_j$ 诱导 $i_j^*:H_T^*(X)\to H_T^*(Z_j)$。此外，$Z_j$ 是 $X$ 的子流形，故还有 Gysin 映射
\[
i_{j!}:H_T^*(Z_j)\longrightarrow H_T^*(X).
\]
这是因为诱导映射 $(Z_j)_T\to X_T$ 可以由子流形的嵌入作有限维近似。Gysin 映射具有如下性质：对任意 $\alpha\in H_T^*(Z_j)$，
\begin{equation}\label{eq:equivariant-self-intersection}
i_j^*\circ i_{j!}(\alpha)=\alpha\cup\operatorname{Euler}_T(N_j).
\end{equation}
这些断言的证明见 \SourceCite[第 6 章附录 A]{Audin}。最后，$H_T^*(Z_j)$ 是 $H^*(BT)\simeq\mathbb C[\lambda_1,\ldots,\lambda_n]$ 模。若
\[
\mathcal R_T\simeq\mathbb C(\lambda_1,\ldots,\lambda_n)
\]
是 $H^*(BT)$ 的分式域，就得到局部化 $H_T^*(Z_j)\otimes_{H^*(BT)}\mathcal R_T$。\SourceCite{AB} 的一个重要观察是，$\operatorname{Euler}_T(N_j)$ 在此局部化中可逆。

在上述设置下，有 \SourceCite{AB} 证明的如下\emph{局部化定理}。
\begin{proposition}\label{prop:atiyah-bott-localization}
存在同构
\[
H_T^*(X)\otimes\mathcal R_T\xrightarrow{\sim}\bigoplus_jH_T^*(Z_j)\otimes\mathcal R_T,
\]
由映射 $\alpha\mapsto(i_j^*(\alpha)/\operatorname{Euler}_T(N_j))_j$ 诱导。其逆映射由 $(\alpha_j)_j\mapsto\sum_ji_{j!}(\alpha_j)$ 诱导。特别地，对任意 $\alpha\in H_T^*(X)\otimes\mathcal R_T$，有
\[
\alpha=\sum_ji_{j!}\left(\frac{i_j^*(\alpha)}{\operatorname{Euler}_T(N_j)}\right).
\]
\end{proposition}
\index{局部化定理|hyperpage}\index{Atiyah--Bott 定理|hyperpage}
看看此定理对 $\mathbb P^r$ 给出什么。
\SourceExampleNumber{9.1.2}{1}
\begin{example}\label{ex:projective-equivariant-localization}
考虑 $T=(\mathbb C^*)^{r+1}$ 在 $X=\mathbb P^r$ 上的作用
\[
(t_0,\ldots,t_r)\cdot(x_0,\ldots,x_r)=(t_0^{-1}x_0,\ldots,t_r^{-1}x_r).
\]
选择逆元，是为了使 $(t_0,\ldots,t_r)$ 在齐次形式 $x_i\in H^0(\mathcal O_{\mathbb P^r}(1))$ 上以乘 $t_i$ 的方式作用。

先定义等变超平面类。$\mathbb P_T^r$ 是 $BT$ 上向量丛 $\mathcal O(-\lambda_0)\oplus\cdots\oplus\mathcal O(-\lambda_r)$ 的射影化。因此，$\mathbb P_T^r=\mathbb P(E_T^*)$，其中 $E^*$ 为例题~\ref{ex:diagonal-equivariant-chern} 中丛 $E$ 的对偶，$\mathbb P$ 表示射影化。这给出重言线丛 $\mathcal O_{\mathbb P_T^r}(1)$，令
\[
p=c_1(\mathcal O_{\mathbb P_T^r}(1))\in H_T^*(\mathbb P^r).
\]
可以验证，$p$ 是等变 Chern 类 $c_1^T(\mathcal O_{\mathbb P^r}(1))$，故称为\emph{等变超平面类}。
\SourcePage{301}
接下来计算 $\mathbb P^r$ 的等变上同调。射影丛上同调的标准公式，例如 \SourceCite[第 606 页]{GH}，给出
\[
(-p)^{r+1}-c_1(E_T^*)(-p)^r+c_2(E_T^*)(-p)^{r-1}+\cdots+(-1)^{r+1}c_{r+1}(E_T^*)=0.
\]
\SourceCite{GH} 中记为 $T$ 的丛，在这里是 $\mathcal O_{\mathbb P_T^r}(-1)$，故 $c_1(T)=-p$。由式~\eqref{eq:equivariant-chern-symmetric}，$c_k(E_T^*)=(-1)^kc_k^T(E)=(-1)^k\sigma_k(\lambda_0,\ldots,\lambda_r)$，因此在 $H_T^*(\mathbb P^r)$ 中有
\begin{equation}\label{eq:projective-equivariant-relation}
\prod_{j=0}^r(p-\lambda_j)=0.
\end{equation}
于是，$\mathbb P^r$ 的等变上同调为
\[
H_T^*(\mathbb P^r)=\mathbb C[p,\lambda_0,\ldots,\lambda_r]\Big/\left(\prod_{j=0}^r(p-\lambda_j)\right).
\]
形式地令 $\lambda_i=0$，再把 $p$ 换成通常超平面类 $H=i_{\mathbb P^r}^*p$，就恢复 $\mathbb P^r$ 通常上同调的描述。

为应用局部化定理，先理解 $T$ 作用的固定点。共有 $r+1$ 个固定点 $q_j$，照常排序，使 $q_j$ 的第 $j$ 个坐标非零，其余为 $0$。注意 $(q_j)_T=\mathbb P(\mathcal O(-\lambda_j))$，所以 $\mathcal O_{\mathbb P_T^r}(-1)$ 在 $(q_j)_T$ 上限制为 $\mathcal O(-\lambda_j)$；相应地，$\mathcal O_{\mathbb P_T^r}(1)$ 限制为 $\mathcal O(\lambda_j)$。因此
\begin{equation}\label{eq:projective-hyperplane-fixed-restriction}\index{等变超平面类|hyperpage}
i_j^*(p)=\lambda_j.
\end{equation}
换句话说，$T$ 在 $\mathbb Cq_j$ 上的作用具有权 $-\lambda_j$，因为按定义 $(t_0,\ldots,t_r)\cdot q_j=t_j^{-1}q_j$。更一般地，$H_T^*(\mathbb P^r)$ 的任何元素都可写为关于 $p$、系数在 $\mathbb C[\lambda_0,\ldots,\lambda_r]$ 中的多项式 $F(p)$，式~\eqref{eq:projective-hyperplane-fixed-restriction} 给出
\begin{equation}\label{eq:projective-polynomial-fixed-restriction}
i_j^*(F(p))=F(\lambda_j).
\end{equation}
因为 $q_j$ 是一个点，法丛 $N_j$ 就是切空间
\[
N_j=T_{q_j}\mathbb P^r=\operatorname{Hom}(\mathbb Cq_j,\mathbb C^{r+1}/\mathbb Cq_j).
\]
所以 $T$ 在 $N_j$ 上的表示具有权 $\lambda_j-\lambda_k$（$k\ne j$），故
\begin{equation}\label{eq:projective-fixed-euler}
\operatorname{Euler}_T(N_j)=\prod_{k\ne j}(\lambda_j-\lambda_k).
\end{equation}
现在计算出命题~\ref{prop:atiyah-bott-localization} 的同构为
\begin{equation}\label{eq:projective-localization-map}
p\longmapsto\left(\frac{\lambda_j}{\prod_{i\ne j}(\lambda_j-\lambda_i)}\right)_{j=0}^r.
\end{equation}
为描述逆映射，令
\begin{equation}\label{eq:projective-fixed-dual-class}
\phi_j=\prod_{k\ne j}(p-\lambda_k).
\end{equation}
于是式~\eqref{eq:projective-polynomial-fixed-restriction} 给出
\begin{equation}\label{eq:projective-dual-class-restriction}
i_k^*(\phi_j)=\begin{cases}\operatorname{Euler}_T(N_j),&k=j,\\0,&k\ne j.\end{cases}
\end{equation}
\SourcePage{302}
因此，命题~\ref{prop:atiyah-bott-localization} 的最后一条断言给出
\[
\phi_j=\sum_{k=0}^ri_{k!}\left(\frac{i_k^*(\phi_j)}{\prod_{l\ne k}(\lambda_k-\lambda_l)}\right)
=i_{j!}\left(\frac{i_j^*(\phi_j)}{\prod_{l\ne j}(\lambda_j-\lambda_l)}\right)=i_{j!}(1).
\]
利用 $H_T^*(q_j)=\mathbb C[\lambda_0,\ldots,\lambda_r]$ 及投影公式 $i_{j!}(i_j^*(\alpha)\cup1)=\alpha\cup i_{j!}(1)$，容易得到式~\eqref{eq:projective-localization-map} 的逆映射为
\[
(\alpha_j)_{j=0}^r\longmapsto\sum_{j=0}^r\alpha_j\cup\phi_j=\sum_{j=0}^r\alpha_j\phi_j.
\]
所以 $\phi_0,\ldots,\phi_r$ 构成 $H_T^*(\mathbb P^r)\otimes\mathcal R_T$ 作为 $\mathcal R_T$ 向量空间的一组基。

最后，取该空间中的任意元素，如上写为关于 $p$、系数在 $\mathcal R_T=\mathbb C(\lambda_0,\ldots,\lambda_r)$ 中的多项式 $F(p)$。再用式~\eqref{eq:projective-polynomial-fixed-restriction}，局部化定理的最后一条断言给出恒等式
\[
F(p)=\sum_{j=0}^rF(\lambda_j)\frac{\prod_{k\ne j}(p-\lambda_k)}{\prod_{k\ne j}(\lambda_j-\lambda_k)}.
\]
这表明，命题~\ref{prop:atiyah-bott-localization} 可以看作 Lagrange 插值公式的广泛推广。
\end{example}
通过向一个点推前，命题~\ref{prop:atiyah-bott-localization} 给出积分公式。对于具有 $T$ 作用的簇 $X$，平凡映射 $X\to\mathrm{point}$ 诱导等变投影 $\pi_X:X_T\to BT$。当 $X$ 完备时，推前 $\pi_{X!}$ 由沿纤维积分给出，在这里称为\emph{等变积分}。因此，常将 $\pi_{X!}$ 写为
\[
\int_{X_T}:H_T^*(X)\longrightarrow H^*(BT).
\]
利用等变积分，得到命题~\ref{prop:atiyah-bott-localization} 的如下推论。
\begin{corollary}\label{cor:localization-integration}
对于任意 $\alpha\in H_T^*(X)\otimes\mathcal R_T$，有
\[
\int_{X_T}\alpha=\sum_j\int_{(Z_j)_T}\frac{i_j^*(\alpha)}{\operatorname{Euler}_T(N_j)}.
\]
\end{corollary}
\index{等变积分|hyperpage}
命题~\ref{prop:atiyah-bott-localization} 和推论~\ref{cor:localization-integration} 给出\emph{局部化}这一概念：将一个类等变地限制到各固定点分支，并除以该分支法丛的等变 Euler 类。这样，一个类的等变积分的计算，就归结为其限制在全部固定点分支上的等变积分之和。

最终要将理论用于轨形 $\overline M_{0,n}(\mathbb P^r,d)$，故实际上需要推论~\ref{cor:localization-integration} 对光滑代数叠的稍加变化的版本。如\hyperref[chap:gromov-witten]{第七章}所述，代数叠的交理论参见 \SourceCite{Vistoli}，代数叠的介绍参见 \SourceCite{BEFFGK}。

设轨形 $X$ 为光滑代数叠的底层簇。即 $X$ 具有\hyperref[app:singular-varieties]{附录 A} 所讨论的局部商图 $U/H$，其中 $U$ 光滑，$H$ 为 $GL(n,\mathbb C)$ 中作用于 $U$ 的小群。在相应光滑代数叠上有 $T$ 作用时，可以在 $T$ 与 $X$ 上局部地将其实现为 $U$ 上的作用。这样，通过局部在 $U$ 中工作，仍能定义固定点轨迹 $Z_j$、包含映射 $i_j:Z_j\hookrightarrow X$ 和等变法丛。因此得到如下代数叠版本。
\SourcePage{303}
\begin{corollary}\label{cor:orbifold-localization-integration}
设 $X$ 为具有 $T$ 作用的光滑代数叠的底层轨形。若 $\alpha\in H_T^*(X)\otimes\mathcal R_T$，则
\[
\int_{X_T}\alpha=\sum_j\int_{(Z_j)_T}\frac{i_j^*(\alpha)}{a_j\operatorname{Euler}_T(N_j)},
\]
其中 $a_j$ 是 $Z_j$ 一般点所对应对象的自同构群的阶。
\end{corollary}

\subsection{Chern 类的多项式}\label{subsec:chern-polynomial-localization}
许多应用中，可更具体地表述推论~\ref{cor:localization-integration}。设具有 $T$ 作用的光滑簇 $X$ 上有等变向量丛 $E$，希望计算 $\int_XP(c_k(E))$，其中 $P$ 是 $E$ 的 Chern 类的多项式。将 $c_k(E)$ 替换为\hyperref[subsec:equivariant-cohomology]{第~\ref*{subsec:equivariant-cohomology} 小节}定义的等变 Chern 类 $c_k^T(E)\in H_T^*(X)$，并沿包含映射 $i_{\mathrm{point}}$ 拉回，就将问题置于等变背景中。为说明这给出 $\int_XP(c_k(E))$，考虑交换图
\begin{equation}\label{eq:equivariant-integral-forgetting-diagram}
\begin{array}{ccc}
X&\longrightarrow&\mathrm{point}\\
\big\downarrow\scriptstyle i_X&&\big\downarrow\scriptstyle i_{\mathrm{point}}\\
X_T&\longrightarrow&BT.
\end{array}
\end{equation}
它意味着 $i_{\mathrm{point}}^*\circ\int_{X_T}=\int_X\circ i_X^*$，故
\[
i_{\mathrm{point}}^*\int_{X_T}P(c_k^T(E))
=\int_Xi_X^*P(c_k^T(E))=\int_XP(c_k(E)),
\]
正如所断言的。关键是 $\int_{X_T}P(c_k^T(E))$ 为可用局部化计算的等变积分。更精确地，将此式与推论~\ref{cor:localization-integration} 结合，就能根据 $E$ 在 $X$ 的 $T$ 固定点集上的行为，计算 $\int_XP(c_k(E))$。

为更明确地说明，考察固定点集 $X^T$ 有限的情形。每个点 $Z_j\in X^T$ 上，限制 $E|_{Z_j}$ 分解为 $T$ 的特征标 $\chi_j^1,\ldots,\chi_j^s$。每个 $\chi_j^k$ 都是已定义的基本特征标 $\rho_1,\ldots,\rho_n$ 的线性组合。这里使用特征标的加法记号，即 $(\rho_1+\rho_2)(t)=\rho_1(t)\rho_2(t)$。将此线性组合写为 $\chi_j^k=\ell_j^k(\rho_1,\ldots,\rho_n)$，简记为 $\chi_j^k=\ell_j^k(\rho)$。

因为 $\rho_i$ 的权为 $\lambda_i$，$\chi_j^k=\ell_j^k(\rho)$ 的权为 $\ell_j^k(\lambda)$，即在线性组合 $\ell_j^k(\rho)$ 中将 $\rho_i$ 替换为 $\lambda_i$ 的结果。因此，$T$ 在 $E|_{Z_j}$ 上的作用具有权 $\ell_j^k(\lambda)$，从而
\[
i_j^*(c_k^T(E))=c_k^T(E|_{Z_j})=\sigma_k(\ell_j^1(\lambda),\ldots,\ell_j^s(\lambda)).
\]
此外，因 $Z_j$ 是点，法丛 $N_j$ 就是切空间 $T_{Z_j}X$。若 $T$ 在该切空间上的特征标为 $t_j^1(\rho),\ldots,t_j^d(\rho)$，其中 $d=\dim X$，则
\begin{equation}\label{eq:isolated-fixed-tangent-euler}
\operatorname{Euler}_T(N_j)=\prod_kt_j^k(\lambda).
\end{equation}
现在对推论~\ref{cor:localization-integration} 的右端施加 $i_{\mathrm{point}}^*$，沿用上述记号，得到如下陈述。
\SourcePage{304}
\begin{proposition}\label{prop:chern-polynomial-localization}
设 $T$ 作用于完备光滑簇 $X$，其固定点 $Z_j$ 孤立，且 $E$ 为 $X$ 上的等变向量丛。若 $P(c_k(E))$ 为 $E$ 的 Chern 类中上同调次数为 $2\dim X$ 的多项式，则
\[
\int_XP(c_k(E))=\sum_j\frac{P\bigl(\sigma_k(\ell_j^1(\lambda),\ldots,\ell_j^s(\lambda))\bigr)}{\prod_kt_j^k(\lambda)},
\]
其中 $\sigma_k$ 为第 $k$ 个初等对称函数。
\end{proposition}
这就是 \SourceCite{ES2} 用于计数应用的表述。若对一般分支 $Z_j$ 重复论证，就得到 \SourceCite{Kontsevich2} 中的表述。

也可将命题~\ref{prop:chern-polynomial-localization} 推广到光滑代数叠：在上述论证中，将推论~\ref{cor:localization-integration} 替换为推论~\ref{cor:orbifold-localization-integration}。结果是，若用轨形描述代数叠，则上述求和中的每一项，都须除以相应 $Z_j$ 对应对象的自同构群的阶。下面\hyperref[subsec:multiple-cover-localization]{第~\ref*{subsec:multiple-cover-localization} 小节}将举例说明。
\SourceExampleNumber{9.1.3}{1}
\begin{example}\label{ex:quintic-lines-localization}
用命题~\ref{prop:chern-polynomial-localization} 计算一般\IndexLocation{indexloc:1182:304}\IndexLocation{indexloc:1188:304}五次三维簇 $V\subset\mathbb P^4$ 上的直线数。例题~\ref{ex:quintic-lines-normal-cone} 曾将 $V$ 上的直线概形实现为 $\mathbb P^4$ 中直线的 Grassmann 簇 $G(2,5)$ 上向量丛 $\operatorname{Sym}^5U^*$ 的一个截面的零点轨迹。$U$ 为 Grassmann 簇上的秩 $2$ 重言子丛，直线 $\ell$ 上的纤维 $U_\ell$ 是射影化为 $\ell$ 的 $\mathbb C^5$ 的二维子空间。$V\subset\mathbb P^4$ 的方程诱导秩 $6$ 丛 $\operatorname{Sym}^5U^*$ 的截面 $s$；例题~\ref{ex:quintic-gw-euler-integral} 表明，直线数是 Euler 类 $c_6(\operatorname{Sym}^5U^*)$ 的次数，即 Gromov--Witten 不变量
\[
\langle I_{0,0,1}\rangle=\int_{G(2,5)}c_6(\operatorname{Sym}^5U^*).
\]
在命题~\ref{prop:chern-polynomial-localization} 中取 $E=\operatorname{Sym}^5U^*$，$P(c_k(E))$ 为其 Euler 类；使用 $(\mathbb C^*)^5$ 在 $\mathbb C^5$ 上的自然作用
\[
(t_1,\ldots,t_5)\cdot(x_1,\ldots,x_5)=(t_1^{-1}x_1,\ldots,t_5^{-1}x_5).
\]
它诱导 $T=(\mathbb C^*)^5$ 在 $G(2,5)$ 上的作用，共有 $10$ 个孤立固定点 $L_I$，对应 $\mathbb P^4$ 的 $10$ 条坐标直线。每个 $L_I$ 由 $I\subset\{1,\ldots,5\}$、$|I|=2$ 标记，并由所有 $j\notin I$ 对应的方程 $x_j=0$ 定义。$L_I$ 的法丛就是 $G(2,5)$ 在 $L_I$ 处的切空间，需要计算它的等变 Euler 类。与式~\eqref{eq:isolated-fixed-tangent-euler} 一样，将切空间的 $T$ 作用分解为特征标，把每个 $\rho_i$ 替换为 $\lambda_i$，然后相乘。

Grassmann 簇的切丛为 $\operatorname{Hom}(U,\mathbb C^5/U)$。关键是，在 $L_I\in G(2,5)$ 处，$U|_{L_I}$ 的特征标为 $-\rho_i$（$i\in I$）。因此，$T$ 在该切空间上的特征标为
\[
\{\rho_i-\rho_j:i\in I,\ j\notin I\},
\]
其等变 Euler 类为
\[
\prod_{i\in I,\,j\notin I}(\lambda_i-\lambda_j).
\]
\SourcePage{305}
另一个所需对象是 $i_I^*(\operatorname{Sym}^5U^*)$ 的等变 Euler 类，其中 $i_I:\{L_I\}\hookrightarrow G(2,5)$ 为包含映射。用 $U^*|_{L_I}$ 的特征标计算；如上，这些特征标为 $\rho_i$（$i\in I$）。因此，写 $I=\{i_1,i_2\}$，得到
\[
\operatorname{Euler}_T\bigl(i_I^*(\operatorname{Sym}^5U^*)\bigr)
=\prod_{a=0}^5\bigl(a\lambda_{i_1}+(5-a)\lambda_{i_2}\bigr).
\]
命题~\ref{prop:chern-polynomial-localization} 现在给出
\[
\int_{G(2,5)}c_6(\operatorname{Sym}^5U^*)
=\sum_{|I|=2}\frac{\prod_{a=0}^5(a\lambda_{i_1}+(5-a)\lambda_{i_2})}{\prod_{i\in I}\prod_{j\notin I}(\lambda_i-\lambda_j)},
\]
仍取 $I=\{i_1,i_2\}$。正如应有的那样，全部 $\lambda_i$ 消去；计算得到答案 $2875$，即一般五次三维簇的直线数。
\end{example}
刚完成的计算涉及将 $\lambda_i$ 的 $10$ 个有理式之和化简；每个有理式的分子、分母都是 $6$ 个关于 $\lambda_i$ 的线性二项式的乘积。手算可能繁琐，计算机却能很快完成。此方法还可迅速进行射影空间中直线、二次曲线和扭曲三次曲线的 Hilbert 概形上的计算，因为这些都是光滑空间，其环面作用也很清楚。例如，完全交 Calabi--Yau 三维簇上的扭曲三次曲线的计算见 \SourceCite{ES2}。这些计算需要两个要素：首先，为所考察的射影曲线找到光滑模空间及 $T$ 作用，使固定点轨迹和权都可计算；其次，将待求数目表示为某个等变向量丛的 Chern 类多项式的次数，而且该丛在固定点处的权可计算。

\section{\texorpdfstring{$\overline M_{0,n}(\mathbb P^r,d)$}{稳定映射模空间}中的局部化}\label{sec:stable-map-localization}
为了求五次三维簇上 $d$ 次有理曲线的数目，Kontsevich 的见解 \SourceCite{Kontsevich2} 是：使用\hyperref[subsec:stable-maps]{第~\ref*{subsec:stable-maps} 小节}描述的稳定映射模空间，就可应用\hyperref[subsec:chern-polynomial-localization]{第~\ref*{subsec:chern-polynomial-localization} 小节}的方法。空间 $\overline M_{0,n}(\mathbb P^r,d)$ 不光滑，但 $\overline{\mathcal M}_{0,n}(\mathbb P^r,d)$ 是光滑代数叠，故推论~\ref{cor:orbifold-localization-integration} 适用。

\subsection{Kontsevich 的方法}\label{subsec:kontsevich-localization}\index{稳定映射!固定点图|hyperpage}
$T=(\mathbb C^*)^{r+1}$ 在 $\mathbb P^r$ 上的自然作用，诱导 $\overline M_{0,n}(\mathbb P^r,d)$ 上的 $T$ 作用。固定点轨迹及其等变法丛已在 \SourceCite{Kontsevich2} 中计算；$g>0$ 的情形，包括虚法丛的局部化公式，见 \SourceCite{GPa}。以下简述主要思想和结果。

$\overline M_{0,n}(\mathbb P^r,d)$ 的 $T$ 固定点由稳定映射 $(f,C,p_1,\ldots,p_n)$ 构成，$C$ 的每个分支 $C_i$ 或被 $f$ 映到 $\mathbb P^r$ 的某个 $T$ 固定点，或多重覆盖一条坐标直线。此外，每个标记点 $p_i$、$C$ 的每个结点、以及 $f$ 的每个分歧点，都映到 $T$ 固定点。因此，若 $C_i$ 通过 $f$ 以次数 $d_i$ 覆盖一条坐标直线，则可选取二者的齐次坐标，使覆盖由 $(x_0,x_1)\mapsto(x_0^{d_i},x_1^{d_i})$ 给出。
\SourcePage{306}
立即可见，该模空间上 $T$ 作用的固定点分支 $Z_j$ 可用组合数据描述。令 $q_0,\ldots,q_r$ 为 $\mathbb P^r$ 的通常 $T$ 固定点。对稳定映射 $(f,C,p_1,\ldots,p_n)$ 关联一棵树 $\Gamma$，其顶点 $v$ 与 $f^{-1}(\{q_0,\ldots,q_r\})$ 的连通分支 $C_v$ 一一对应。每个 $C_v$ 或为 $C$ 中一个点，或为其不可约分支的连通并。$\Gamma$ 的边 $e$ 对应 $C$ 中被 $f$ 映满某条坐标直线 $\ell_e\subset\mathbb P^r$ 的不可约分支 $C_e$。

树 $\Gamma$ 还带有各种标记。对每个顶点 $v$，关联由 $f(C_v)=q_{i_v}$ 定义的数 $i_v$，以及由所有满足 $p_i\in C_v$ 的 $i$ 组成的集合 $S_v$。对每条边 $e$，关联映射 $f|_{C_e}:C_e\to\ell_e$ 的次数 $d_e$。于是，固定点集 $\overline M_{0,n}(\mathbb P^r,d)^T$ 的连通分支，与带有标记 $i_v,S_v,d_e$ 并满足以下条件的连通树一一对应：
\begin{itemize}
\item 若边 $e$ 含顶点 $v,v'$，则 $i_v\ne i_{v'}$，且 $\ell_e$ 为连接 $q_{i_v}$ 与 $q_{i_{v'}}$ 的坐标直线。
\item $\sum_ed_e=d$。
\item $\{1,\ldots,n\}=\coprod_vS_v$。
\end{itemize}
称数据 $(\Gamma,i_v,S_v,d_e)$ 为稳定映射的\emph{图}，通常将全部数据简记为 $\Gamma$。回顾，顶点 $v$ 的\emph{价数} $\operatorname{val}(v)$ 是与 $v$ 相接的边数。\index{稳定映射!固定点图|hyperpage}\index{价数|hyperpage}

具有固定图 $\Gamma$ 的稳定映射，自然定义子叠
\[
\overline{\mathcal M}_\Gamma\subset\overline{\mathcal M}_{0,n}(\mathbb P^r,d).
\]
为研究它，取 $(f,C,p_1,\ldots,p_n)\in\overline{\mathcal M}_\Gamma$。对于 $C_v$ 是曲线的顶点 $v$，$C_v$ 有 $n(v)=|S_v|+\operatorname{val}(v)$ 个特殊点：其上有 $|S_v|$ 个标记点，另有 $\operatorname{val}(v)$ 个结点，使 $C_v$ 与包含 $v$ 的边 $e$ 对应的不可约分支 $C_e$ 相交。实际上，$C_v$ 连同这些 $n(v)$ 个点组成稳定曲线，给出 $\overline M_{0,n(v)}$ 的元素。

更仔细的描述如下。给定一组 $n(v)$ 点亏格零稳定曲线 $\{C_v\}\in\prod_{\dim C_v=1}\overline M_{0,n(v)}$，每当边 $e$ 含顶点 $v,v'$，就用一条新曲线 $C_e\simeq\mathbb P^1$ 的 $0,\infty$ 将 $C_v$ 与 $C_{v'}$ 连接，从而构造曲线 $C$。这可以系统地完成。再定义 $f:C\to\mathbb P^r$，将每个曲线 $C_v$ 收缩到 $q_{i_v}$，并令 $f|_{C_e}$ 为 $(x_0,x_1)\mapsto(x_0^{d_e},x_1^{d_e})$。得到态射
\[
\psi_\Gamma:\prod_{\dim C_v=1}\overline M_{0,n(v)}\longrightarrow\overline{\mathcal M}_\Gamma.
\]
$n(v)$ 对所有顶点都有定义，而且 $\dim C_v=1$，当且仅当 $C_v$ 含有被 $f$ 收缩的分支，也当且仅当 $n(v)\ge3$。

将上述乘积定义为 $\overline M_\Gamma$；若没有收缩分支，就令 $\overline M_\Gamma$ 为一个点。上述 $\psi_\Gamma:\overline M_\Gamma\to\overline{\mathcal M}_\Gamma$ 是有限态射，并非同构。用代数叠语言说，有有限自同构群 $A_\Gamma$ 作用于 $\overline M_\Gamma$，其代数叠意义下的商为 $\overline{\mathcal M}_\Gamma$。群 $A_\Gamma$ 处于正合列
\begin{equation}\label{eq:fixed-graph-automorphisms}\index{稳定映射!自同构群|hyperpage}
1\longrightarrow\prod_e\mathbb Z/d_e\mathbb Z\longrightarrow A_\Gamma
\longrightarrow\operatorname{Aut}(\Gamma)\longrightarrow1,
\end{equation}
其中 $\operatorname{Aut}(\Gamma)$ 为保持标记的图自同构群。此列可解释如下：子群 $\prod_e\mathbb Z/d_e\mathbb Z$ 来自 $d_e$ 重覆盖的覆盖
\SourcePage{307}
变换。商映射如下产生：每个 $g\in A_\Gamma$ 将一组 $n(v)$ 点亏格零稳定曲线 $\{C_v\}\in\overline M_\Gamma$ 映到另一组 $\{C'_v\}$，使 $f=\psi_\Gamma(\{C_v\})$ 与 $f'=\psi_\Gamma(\{C'_v\})$ 是同构的稳定映射。该同构将 $f^{-1}(q_i)$ 与 $(f')^{-1}(q_i)$ 的连通分支等同。特别地，每个 $C_v$ 与某个 $C'_w$ 等同，且 $i_v=i_w$。容易看出，顶点集的置换 $v\mapsto w$ 与 $C_v$ 的选择无关，并诱导保持标记的 $\Gamma$ 自同构 $\alpha(g)$。映射 $g\mapsto\alpha(g)$ 定义式~\eqref{eq:fixed-graph-automorphisms} 中的 $A_\Gamma\to\operatorname{Aut}(\Gamma)$，该列的正合性随即可见。

因此，将推论~\ref{cor:orbifold-localization-integration} 应用于 $\overline M_{0,n}(\mathbb P^r,d)$，此时与固定点分支 $\overline M_\Gamma$ 对应的项，相应项的分母因子 $a_\Gamma$ 就是 $A_\Gamma$ 的阶。
\SourceExampleNumber{9.2.1}{1}
\begin{example}\label{ex:fixed-graphs-p1-degree2}
计算标准 $T$ 作用在 $\overline M_{0,0}(\mathbb P^1,2)$ 上的固定点分支。可能的图 $\Gamma$ 为：
\begin{figure}[H]\centering
\includegraphics[width=.785\textwidth]{fixed-graphs-p1-degree2.png}
\end{figure}
图中 $d_e$ 放在边 $e$ 上方，$i_v$ 放在顶点 $v$ 下方。没有标记点，故全部 $S_v$ 为空。

得到固定点集的三个分支，相应 $T$ 固定稳定映射分别如下：
\begin{enumerate}[label=(\roman*)]
\item 曲线 $C=C_1\cup C_2$ 有两个分支。$f$ 在每个分支上的限制 $f|_{C_i}:C_i\xrightarrow{\sim}\mathbb P^1$ 为同构，且 $f(C_1\cap C_2)=q_1$。
\item 曲线 $C=C_1\cup C_2$ 有两个分支。$f$ 在每个分支上的限制 $f|_{C_i}:C_i\xrightarrow{\sim}\mathbb P^1$ 为同构，且 $f(C_1\cap C_2)=q_0$。
\item 曲线 $C$ 不可约，可在 $C\simeq\mathbb P^1$ 上取坐标 $(z_0,z_1)$，使 $f(z_0,z_1)=(z_0^2,z_1^2)$。
\end{enumerate}
每种情形都有 $A_\Gamma\simeq\mathbb Z_2$，故 $a_\Gamma=2$。在 (i)、(ii) 中，非平凡自同构来自 $\Gamma$ 的显然自同构；在 (iii) 中，它来自二重覆盖。
\end{example}
\SourceExampleNumber{9.2.1}{2}
\begin{example}\label{ex:fixed-graphs-p2-marked-degree2}
计算标准 $T$ 作用在 $\overline M_{0,1}(\mathbb P^2,2)$ 上的固定点分支。$\mathbb P^2$ 的 $T$ 作用有三个固定点
\[
q_0=(1,0,0),\qquad q_1=(0,1,0),\qquad q_2=(0,0,1).
\]
相应的可能图为：
\begin{figure}[H]\centering
\includegraphics[width=.792\textwidth]{fixed-graphs-p2-marked-degree2.png}
\end{figure}
只要 $S_v$ 非空，就在顶点 $v$ 上方加入标记 $S_v$。

第二个图对应的稳定映射，其源曲线为 $C=C_1\cup C_2\cup C_3$，$C_2$ 分别与 $C_1,C_3$ 相交于一个点，没有其他交点。$f$ 将 $C_1$ 同构地映到连接 $q_i,q_j$ 的直线，将 $C_2$ 映到点 $q_j$，并将 $C_3$
\SourcePage{308}
同构地映到连接 $q_j,q_k$ 的直线。标记点 $p_1$ 是 $C_2$ 的光滑点。$f$ 稳定，因为 $C_2$ 含三个特殊点 $p_1,C_1\cap C_2,C_2\cap C_3$。从图的中间顶点 $v$ 可立即识别收缩分支的存在，因为 $n(v)=1+2=3$。
\end{example}
为在 $\overline M_{0,n}(\mathbb P^r,d)$ 中开展计算，使用推论~\ref{cor:orbifold-localization-integration}，故须描述固定点分支的法丛。将 $\overline{\mathcal M}_\Gamma$ 的法丛记为 $N_\Gamma$。其等变 Euler 类计算的参考文献包括 \SourceCite{Kontsevich2,Givental4,GPa}；其中 \SourceCite{GPa} 还推广到 $\overline M_{g,n}(\mathbb P^r,d)$，需要等变虚拟基本类的\IndexLocation{indexloc:1417:308}局部化公式。

回顾\hyperref[subsec:virtual-fundamental-class]{第~\ref*{subsec:virtual-fundamental-class} 小节}，代数叠 $\overline{\mathcal M}_{0,n}(\mathbb P^r,d)$ 在稳定映射 $(f,C,p_1,\ldots,p_n)$ 处的切空间是超 Ext 群
\[
\operatorname{Ext}^1_C(f^*\Omega^1_{\mathbb P^r}\longrightarrow\Omega_C^1(D),\mathcal O_C),
\]
其中 $D=\sum_ip_i$。用 Ext 的长正合列理解此群，得到
\begin{equation}\label{eq:stable-map-tangent-exact-localization}
\begin{aligned}
0\longrightarrow\operatorname{Hom}(\Omega_C^1(D),\mathcal O_C)
&\longrightarrow H^0(C,f^*T_{\mathbb P^r})\\
&\longrightarrow\operatorname{Ext}^1(f^*\Omega^1_{\mathbb P^r}\to\Omega_C^1(D),\mathcal O_C)\\
&\longrightarrow\operatorname{Ext}^1(\Omega_C^1(D),\mathcal O_C)\longrightarrow0.
\end{aligned}
\end{equation}
同时，$\operatorname{Ext}^2(f^*\Omega^1_{\mathbb P^r}\to\Omega_C^1(D),\mathcal O_C)=0$，因为 $\mathbb P^r$ 凸，故 $H^1(C,f^*T_{\mathbb P^r})=0$。对更一般的 $g,X$，须用计算切空间—\IndexLocation{indexloc:1328:308}障碍复形的类似正合列替换式~\eqref{eq:stable-map-tangent-exact-localization}，如\hyperref[subsec:virtual-fundamental-class]{第~\ref*{subsec:virtual-fundamental-class} 小节}所述。

在 $\overline{\mathcal M}_\Gamma\subset\overline M_{0,n}(\mathbb P^r,d)$ 的 $T$ 固定点处，可由此正合列求切空间的 $T$ 作用的权。再考察各项的 $T$ 作用，并限制到“移动部分”，就得到法丛 $N_\Gamma$ 的权。将这些权相乘，照常得到其等变 Euler 类。

给出 $\operatorname{Euler}_T(N_\Gamma)$ 的公式前，需要一些记号。按 \SourceCite{Kontsevich2}，\emph{旗} $F$ 是一对 $(v,e)$，其中 $e$ 为含顶点 $v$ 的边。令 $i(F)=i_v$，$j(F)=i_{v'}$，其中 $v'$ 为 $e$ 的另一个顶点。对旗 $F=(v,e)$，令
\[
\omega_F=\frac{\lambda_{i(F)}-\lambda_{j(F)}}{d_e}.
\]
它是 $H^2(BT)$ 中对应于如下权的元素：$T$ 在边 $e$ 对应的分支 $C_e$ 上、位于 $q_{i_v}$ 上方的点 $p_F$ 处切空间上的作用的权。当 $n(v)\ge3$ 时，令 $e_F$ 为 $\overline M_\Gamma$ 上某线丛的第一 Chern 类，其纤维为顶点 $v$ 对应曲线分支在 $p_F$ 处的余切空间。

若 $\operatorname{val}(v)=1$，将含 $v$ 的唯一旗记为 $F(v)$。若 $\operatorname{val}(v)=2$，将两面旗记为 $F_1(v),F_2(v)$。\index{旗!稳定映射图|hyperpage}
\begin{theorem}\label{thm:fixed-graph-normal-euler}\index{等变 Euler 类!法丛|hyperpage}
法丛 $N_\Gamma$ 的等变 Euler 类，是旗、顶点、边的贡献之积。更精确地，
\[
\operatorname{Euler}_T(N_\Gamma)=e_\Gamma^{\mathrm F}e_\Gamma^{\mathrm v}e_\Gamma^{\mathrm e},
\]
\SourcePage{309}
其中三者定义为
\[
\begin{aligned}
e_\Gamma^{\mathrm F}&=\frac{\prod_{F=(v,e),\,n(v)\ge3}(\omega_F-e_F)}
{\prod_{F=(v,e)}\prod_{j\ne i(F)}(\lambda_{i(F)}-\lambda_j)},\\
e_\Gamma^{\mathrm v}&=\prod_v\prod_{j\ne i_v}(\lambda_{i_v}-\lambda_j)
\frac{\prod_{\operatorname{val}(v)=2,\,S_v=\varnothing}(\omega_{F_1(v)}+\omega_{F_2(v)})}
{\prod_{\operatorname{val}(v)=1,\,S_v=\varnothing}\omega_{F(v)}},\\
e_\Gamma^{\mathrm e}&=\prod_{e=(v,v')}\left[
\frac{(-1)^{d_e}(d_e!)^2(\lambda_i-\lambda_j)^{2d_e}}{d_e^{2d_e}}
\prod_{\substack{a+b=d_e,\ a,b\ge0\\k\ne i,j}}
\left(\frac{a\lambda_i+b\lambda_j}{d_e}-\lambda_k\right)\right],
\end{aligned}
\]
边项中 $i=i_v$、$j=i_{v'}$。
\end{theorem}
\begin{proof}
完整证明见 \SourceCite{Kontsevich2,GPa}。这里只简述涉及下面引理~\ref{lem:node-smoothing-tangent-tensor} 的一小部分证明，此引理以后还会有用。

先注意，$\operatorname{Hom}(\Omega_C^1(D),\mathcal O_C)$ 是 $(C,p_1,\ldots,p_n)$ 的无穷小自同构空间，而 $H^0(C,f^*T_{\mathbb P^r})$ 是映射 $f$ 的一阶形变空间。因此，
\[
H^0(C,f^*T_{\mathbb P^r})/\operatorname{Hom}(\Omega_C^1(D),\mathcal O_C)
\]
是整个稳定映射模空间在 $f$ 处切空间的子空间 $S_{\Gamma,f}$，对应于固定带标记源曲线、仅改变映射的形变。结合式~\eqref{eq:stable-map-tangent-exact-localization}，得到正合列
\[
0\longrightarrow S_{\Gamma,f}\longrightarrow T_f\overline M_{0,n}(\mathbb P^r,d)
\xrightarrow{\kappa}\operatorname{Ext}^1(\Omega_C^1(D),\mathcal O_C)\longrightarrow0.
\]
此外，由 Ext 的局部到整体谱序列，还有映射
\[
\mu:\operatorname{Ext}^1(\Omega_C^1(D),\mathcal O_C)
\longrightarrow H^0(C,\mathcal E xt^1(\Omega_C^1(D),\mathcal O_C)).
\]
$C$ 是曲线，谱序列在 $E_2$ 退化，故 $\mu$ 满射。$\mathcal E xt^1(\Omega_C^1(D),\mathcal O_C)$ 是支集位于 $C$ 的结点的挠层，因为 $\Omega_C^1(D)$ 在光滑点处局部自由。立即得到
\[
\mathcal E xt^1(\Omega_C^1(D),\mathcal O_C)\simeq\mathcal E xt^1(\Omega_C^1,\mathcal O_C),
\]
因而可将 $\mu$ 看作满射
\[
\mu:\operatorname{Ext}^1(\Omega_C^1(D),\mathcal O_C)
\longrightarrow H^0(C,\mathcal E xt^1(\Omega_C^1,\mathcal O_C)).
\]
由 \SourceCite{DM}，右端是忘掉 $f$ 和全部标记点后，曲线 $C$ 的一阶光滑化空间。尽管此处 $C$ 未必稳定，\SourceCite{DM} 的论证仍然适用。

这说明，$\kappa^{-1}(\mu^{-1}(\{0\}))\subset T_f\overline M_{0,n}(\mathbb P^r,d)$ 是保持全部结点的一阶形变空间。由于来自 $T_f\overline M_\Gamma$ 的形变显然保持这些结点，因此 $T_f\overline M_\Gamma\subset\kappa^{-1}(\mu^{-1}(\{0\}))$。于是 $\mu\circ\kappa$ 诱导满射
\[
N_{\Gamma,f}\longrightarrow H^0(C,\mathcal E xt^1(\Omega_C^1,\mathcal O_C)),
\]
它与右端的自然 $T$ 作用相容。

现给出 \SourceCite{Kontsevich2} 的一个重要引理，显式计算此空间。令 $q_i$ 为 $C$ 的结点，$C_i^1,C_i^2$ 为含该点的两个分支。
\SourcePage{310}
\begin{lemma}\label{lem:node-smoothing-tangent-tensor}\index{结点!光滑化|hyperpage}
存在与 $T$ 作用相容的自然同构
\[
H^0(C,\mathcal E xt^1(\Omega_C^1,\mathcal O_C))
\simeq\bigoplus_i\bigl(T_{q_i}C_i^1\otimes T_{q_i}C_i^2\bigr).
\]
\end{lemma}
\begin{proof}
计算纯属局部，故令 $C$ 为局部曲面 $S$ 内由 $xy=0$ 定义的局部解析曲线，$(x,y)$ 为 $S$ 的局部坐标。将 $C$ 的分支记为 $C_1,C_2$，令 $q=(0,0)$。用标准正合列计算 $\mathcal E xt^1$：
\begin{equation}\label{eq:node-conormal-sequence}
0\longrightarrow I_C/I_C^2\longrightarrow\Omega_S^1|_C\longrightarrow\Omega_C^1\longrightarrow0,
\end{equation}
其中第一个映射将 $[f]$ 映到 $df|_C$，第二个映射是到 $C$ 的拉回。施加 $\mathcal E xt^*(\,\cdot\,,\mathcal O_C)$ 的长正合列，得到
\begin{equation}\label{eq:node-ext-cokernel}
\mathcal E xt^1(\Omega_C^1,\mathcal O_C)
\simeq\operatorname{coker}\left(T_S|_C\xrightarrow h(I_C/I_C^2)^*\right),
\end{equation}
其中 $(I_C/I_C^2)^*=\mathcal H om(I_C/I_C^2,\mathcal O_C)$。计算可知，$h(\partial/\partial x)$ 将 $I_C/I_C^2$ 的生成元 $xy$ 映到 $y$，而 $h(\partial/\partial y)$ 将 $xy$ 映到 $x$。

内在地，这给出同构
\begin{equation}\label{eq:node-local-tensor-isomorphism}
T_qC_1\otimes T_qC_2\simeq\mathcal E xt^1(\Omega_C^1,\mathcal O_C)_q,
\end{equation}
借助式~\eqref{eq:node-ext-cokernel} 的等同，定义为
\begin{equation}\label{eq:node-local-tensor-map}
v_1\otimes v_2\longmapsto[f\longmapsto v_1v_2(f)(q)].
\end{equation}
为理解此式，将 $v_1,v_2$ 看作 $S$ 上向量场在 $C_1,C_2$ 上的限制，且每个 $v_i$ 与 $C_i$ 相切。将所得函数 $v_1v_2f$ 在 $q$ 处取值。给定向量场 $v_1,v_2$，映射 $f\mapsto v_1v_2(f)(q)$ 给出 $(I_C/I_C^2)^*|_q$ 的元素，并通过式~\eqref{eq:node-ext-cokernel} 确定余核中的类。容易看出，这个类仅依赖点 $q$ 处的 $v_1\otimes v_2$。用局部坐标 $x,y$ 表示，式~\eqref{eq:node-local-tensor-map} 将生成元 $(\partial/\partial x)\otimes(\partial/\partial y)$ 映到把 $xy$ 送到 $1$ 的 $(I_C/I_C^2)^*$ 元素的类。因此，它确实定义了同构~\eqref{eq:node-local-tensor-isomorphism}。

该同构的自然性说明，它与两侧的自然 $T$ 作用相容。
\end{proof}
在定理~\ref{thm:fixed-graph-normal-euler} 的证明中，引理~\ref{lem:node-smoothing-tangent-tensor} 在 $\operatorname{Euler}_T(N_\Gamma)$ 的公式中用于两处：
\begin{itemize}
\item $e_\Gamma^{\mathrm F}$ 中的 $\omega_F-e_F$ 项，来自收缩分支与以正次数映到 $\mathbb P^r$ 的分支相遇的结点。
\item $e_\Gamma^{\mathrm v}$ 中的 $\omega_{F_1(v)}+\omega_{F_2(v)}$ 项，来自两个都以正次数映到 $\mathbb P^r$ 的分支相遇的结点。
\end{itemize}
为完成证明，还须理解 $e_\Gamma^{\mathrm F},e_\Gamma^{\mathrm v},e_\Gamma^{\mathrm e}$ 的其他项。例如，可证明：
\begin{itemize}
\item $e_\Gamma^{\mathrm v}$ 的分母来自 $T$ 在 $\operatorname{Hom}(\Omega_C^1,\mathcal O_C)$ 中无穷小自同构上的作用的权。
\item $e_\Gamma^{\mathrm F}$ 的分母，以及 $e_\Gamma^{\mathrm v}$ 的第一个因子，分别来自 $\mathbb P^r$ 在 $q_{i(F)}$、$q_{i_v}$ 处切空间的贡献。
\end{itemize}
这些由式~\eqref{eq:stable-map-tangent-exact-localization} 不难得到。不过，还没有讨论 $e_\Gamma^{\mathrm e}$ 中的项，所以要证明定理，仍有许多工作。如前所述，完整细节见 \SourceCite{Kontsevich2,GPa}。
\end{proof}

\SourcePage{311}
定理~\ref{thm:fixed-graph-normal-euler} 给出的 $\operatorname{Euler}_T(N_\Gamma)$ 公式中，三个因子 $e_\Gamma^{\mathrm F},e_\Gamma^{\mathrm v},e_\Gamma^{\mathrm e}$ 之间显然有许多约去的项。因此，在具体例题中，这个公式未必是计算效率最高的形式。定理的优点，是给出了容易自动化的 $\operatorname{Euler}_T(N_\Gamma)$ 计算算法。下面的第~\ref{subsec:multiple-cover-localization} 小节将详细计算一个非平凡例题，并验证这些公式的正确性。

应当指出，引理~\ref{lem:node-smoothing-tangent-tensor} 中的等同
\[
H^0(C,\mathcal E xt^1(\Omega_C^1,\mathcal O_C))
\simeq\bigoplus_i(T_{q_i}C_i^1\otimes T_{q_i}C_i^2)
\]
已隐含地出现在 \SourceCite{Mumford2,HM} 中。下面第~\ref{subsec:multiple-cover-localization} 小节还将使用这一引理。\hyperref[chap:mirror-theorem]{第~\ref*{chap:mirror-theorem} 章}所述镜像定理的证明中，它也起着核心作用。

在更一般的空间 $\overline M_{0,n}(X,\beta)$ 上计算时，引理右侧对应的 $T$ 表示仍自然出现，尽管引理按原有表述不能直接适用。这源于切空间与虚拟基本类之间的联系；二者都由 $\operatorname{Ext}^*(f^*\Omega_X^1\to\Omega_C^1(D),\mathcal O_C)$ 的长正合列产生。

\SourceExampleNumber{9.2.1}{3}
\begin{example}\label{ex:quintic-quartics-localization}\index{五次三维簇!有理四次曲线|hyperpage}
这里建立的理论最早的应用之一，是计算五次三维簇的 Gromov--Witten 不变量 $N_4=\langle I_{0,0,4}\rangle$ \SourceCite{Kontsevich2}。回忆例题~\ref{ex:quintic-gw-euler-integral} 中的公式，
\[
N_4=\int_{\overline M_{0,0}(\mathbb P^4,4)}\operatorname{Euler}(\mathcal V_4),
\]
其中 $\mathcal V_4$ 是 $\overline M_{0,0}(\mathbb P^4,4)$ 上的向量丛，在 $(f,C)$ 处的纤维为 $H^0(C,f^*\mathcal O_{\mathbb P^4}(5))$。

对于 $T=(\mathbb C^*)^5$ 在 $\mathbb P^4$ 上的作用，可以列出所有可能的图 $\overline M_\Gamma$，并计算相应的等变 Euler 类 $\operatorname{Euler}_T(N_\Gamma)$。还可以将上同调类显式表示成 Čech 上闭链，计算 $\operatorname{Euler}(\mathcal V_4)$ 的限制的权。第~\ref{subsec:multiple-cover-localization} 小节将解释这一做法。再施加命题~\ref{prop:chern-polynomial-localization}，所得计算规模很大；不过，正如 \SourceCite{Kontsevich2} 所说，“在 Sun 上计算 5 分钟”便得到
\[
N_4=\frac{15517926796875}{64}.
\]
用瞬子数表示，有 $N_4=n_1\,4^{-3}+n_2\,2^{-3}+n_4$。代入已知值 $n_1=2875$ 和 $n_2=609250$，Kontsevich 得到
\[
n_4=242467530000.
\]
这是首次严格计算出 $n_4$，证实了镜像对称的预言 \SourceCite{CdGP}。正如例题~\ref{ex:quintic-instanton-enumerative} 所说明的，这意味着一般的五次三维簇恰有 $242467530000$ 条有理四次曲线。
\end{example}

\subsection{多重覆盖与 Gromov--Witten 不变量}\label{subsec:multiple-cover-localization}
设 $V$ 是 Calabi--Yau 三维簇，$C\subset V$ 是光滑的无穷小刚性有理曲线。掌握局部化以后，我们便可以介绍 Kontsevich 的方法，计算 $C$ 的 $d$ 次覆盖对 Gromov--Witten 不变量 $\langle I_{0,0,d[C]}\rangle$ 的贡献。根据定理~\ref{thm:multiple-cover}，答案是 $1/d^3$。\index{多重覆盖|hyperpage}

固定同构 $\mathbb P^1\simeq C$。刚性假设意味着 $C\subset V$ 的法丛为 $N\simeq\mathcal O_{\mathbb P^1}(-1)\oplus\mathcal O_{\mathbb P^1}(-1)$。
\SourcePage{312}
虚拟基本类 $[\overline M_{0,0}(V,d[C])]^{\mathrm{virt}}$ 是一个零维圈；按定义，$\langle I_{0,0,d[C]}\rangle$ 是它的次数。我们要计算其中支集位于 $C$ 的 $d$ 次多重覆盖簇 $M_{d,C}(V)$ 上的那一部分的贡献。在定理~\ref{thm:multiple-cover} 前的讨论中，已证明 $M_{d,C}(V)$ 是 $\overline M_{0,0}(V,d[C])$ 的一个连通分支。

同构 $\mathbb P^1\simeq C$ 诱导同构 $M_{d,C}(V)\simeq\overline M_{0,0}(\mathbb P^1,d)$，计算就在后一个空间上进行。容易看出，$\dim\overline M_{0,0}(\mathbb P^1,d)=2d-2$。证明的第一步，是将 $[\overline M_{0,0}(V,d[C])]^{\mathrm{virt}}$ 在该空间上的限制，表示为适当向量丛的 Euler 类。

具体如下。考虑稳定映射 $f\in\overline M_{0,0}(\mathbb P^1,d)$。拉回定义 $C$ 在 $V$ 中法丛的正合列，得到
\[
0\longrightarrow f^*T_{\mathbb P^1}\longrightarrow f^*T_V
\longrightarrow f^*(\mathcal O_{\mathbb P^1}(-1)\oplus\mathcal O_{\mathbb P^1}(-1))\longrightarrow0.
\]
命题~\ref{prop:excess-obstruction-euler-class} 的假设成立，因为
\[
h^1(C,f^*T_V)=h^1(C,f^*(\mathcal O_{\mathbb P^1}(-1)\oplus\mathcal O_{\mathbb P^1}(-1)))=2d-2,
\]
最后一个等式来自 Riemann--Roch 定理。因此，该命题说明，虚拟基本类的限制就是秩为 $2d-2$ 的障碍丛 $R^1\pi_{1*}e_1^*(\mathcal O_{\mathbb P^1}(-1)\oplus\mathcal O_{\mathbb P^1}(-1))$ 的 Euler 类。这里的万有稳定映射由曲线族 $\pi_1:\overline M_{0,1}(\mathbb P^1,d)\to\overline M_{0,0}(\mathbb P^1,d)$ 及万有映射 $e_1:\overline M_{0,1}(\mathbb P^1,d)\to\mathbb P^1$ 给出。

于是，\hyperref[chap:gromov-witten]{第七章}的定理~\ref{thm:multiple-cover} 可以改述如下。
\begin{theorem}\label{thm:multiple-cover-localization}\index{多重覆盖公式|hyperpage}
设 $C\subset V$ 是刚性嵌入 Calabi--Yau 三维簇的光滑有理曲线。则 $C$ 的 $d$ 次多重覆盖对 Gromov--Witten 不变量 $\langle I_{0,0,d[C]}\rangle$ 的贡献为
\[
\int_{\overline M_{0,0}(\mathbb P^1,d)}
\operatorname{Euler}\bigl(R^1\pi_{1*}e_1^*(\mathcal O_{\mathbb P^1}(-1)\oplus\mathcal O_{\mathbb P^1}(-1))\bigr)=\frac1{d^3}.
\]
\end{theorem}
\begin{proof}
证明的思路，是对 $T=(\mathbb C^*)^2$ 在 $\mathbb P^1$ 上的通常作用，使用代数叠的局部化公式，即推论~\ref{cor:orbifold-localization-integration}，计算积分。

先考虑 $d=2$。例题~\ref{ex:fixed-graphs-p1-degree2} 已确定 $T$ 在 $\overline M_{0,0}(\mathbb P^1,2)$ 上作用的三个固定点分支。将该例中的图依次记为 $\Gamma_1,\Gamma_2,\Gamma_3$。对于每个图 $\Gamma_i$，计算 $T$ 在法丛 $N_{\Gamma_i}$ 上，以及在障碍丛 $R^1\pi_{1*}e_1^*(\mathcal O_{\mathbb P^1}(-1)\oplus\mathcal O_{\mathbb P^1}(-1))$ 限制到 $Z_{\Gamma_i}$ 上的作用的权。

从 $\Gamma_1$ 开始。它是下面这个带标号的图：
\begin{figure}[H]
\centering\includegraphics[width=.171\textwidth]{multiple-cover-graph-chain.png}
\end{figure}
相应的固定点分支 $Z_{\Gamma_1}$ 只有一点，对应稳定映射 $f_1:C_1\cup C_2\to\mathbb P^1$，其中 $f_1|_{C_i}$ 是同构，且 $f_1(p)=q_1$，$\{p\}=C_1\cap C_2$。

先计算等变 Euler 类。图中有 4 个旗、3 个顶点和 2 条边。利用定理~\ref{thm:fixed-graph-normal-euler} 的公式，得到
\begin{align*}
e_{\Gamma_1}^{\mathrm F}&=\frac1{(\lambda_0-\lambda_1)^2(\lambda_1-\lambda_0)^2},\\
e_{\Gamma_1}^{\mathrm v}&=((\lambda_0-\lambda_1)^2(\lambda_1-\lambda_0))
\frac{2\lambda_1-2\lambda_0}{(\lambda_0-\lambda_1)^2},\\
e_{\Gamma_1}^{\mathrm e}&=(\lambda_0-\lambda_1)^4.
\end{align*}
\SourcePage{313}
立即得出
\begin{equation}\label{eq:degree-two-chain-normal-euler}
\operatorname{Euler}_T(N_{\Gamma_1})=2(\lambda_1-\lambda_0)^2.
\end{equation}
为进一步理解这些公式的来源，下面用更初等的方法计算切空间的权，验证此式。

为描述 $\overline M_{0,0}(\mathbb P^1,2)$ 在 $f_1$ 处的二维切空间，考虑子叠 $\overline M(\mathbb P^1,2)_1\subset\overline M_{0,0}(\mathbb P^1,2)$，由零标记稳定映射 $f:C_1\cup C_2\to\mathbb P^1$ 构成，其中 $f|_{C_i}$ 的次数为 $1$。注意 $f_1\in\overline M(\mathbb P^1,2)_1$。存在态射 $\overline M(\mathbb P^1,2)_1\to\mathbb P^1$，将 $f$ 映到 $f(C_1\cap C_2)$。由于存在交换 $C_1,C_2$ 的自同构，该态射作为代数叠的映射，其次数为 $2$。固定点 $Z_{\Gamma_1}=\{f_1\}$ 由此映到 $q_1$。该自同构对应图 $\Gamma_1$ 的一个自同构，在施加局部化公式时，会在分母中额外贡献因子 $2$。

子叠在 $f_1$ 处的切空间，给出 $T_{f_1}\overline M_{0,0}(\mathbb P^1,2)$ 的一个切方向，所做的只是移动点 $q_1$。因此，$T$ 在这个方向上的作用的权，等于它在 $\mathbb P^1$ 于 $q_1$ 处切空间上的权。用通常方法容易算得此权为 $\lambda_1-\lambda_0$。

为求剩下的权，对定理~\ref{thm:fixed-graph-normal-euler} 的证明中用到的映射作适当的图表追踪，可知商空间
\[
T_{f_1}\overline M_{0,0}(\mathbb P^1,2)/T_{f_1}\overline M(\mathbb P^1,2)_1
\]
自然同构于
\[
H^0(C_1\cup C_2,\mathcal E xt^1(\Omega_{C_1\cup C_2}^1,\mathcal O_{C_1\cup C_2})).
\]
引理~\ref{lem:node-smoothing-tangent-tensor} 又说明，这个商同构于 $T_pC_1\otimes T_pC_2$。

所有这些同构都保持自然的 $T$ 作用。而每个 $T_pC_i$ 作为 $T$ 表示都同构于 $T_{q_1}\mathbb P^1$，所以剩下的权为 $2(\lambda_1-\lambda_0)$。于是切空间的两个权是 $\lambda_1-\lambda_0$ 和 $2(\lambda_1-\lambda_0)$，等变 Euler 类为 $2(\lambda_1-\lambda_0)^2$，与式~\eqref{eq:degree-two-chain-normal-euler} 一致。

现在考虑障碍丛 $R^1\pi_{1*}e_1^*(\mathcal O_{\mathbb P^1}(-1)\oplus\mathcal O_{\mathbb P^1}(-1))$。计算其 $T$ 权，只须求 $R^1\pi_{1*}e_1^*\mathcal O_{\mathbb P^1}(-1)$ 的权，再重复一次答案。$f_1^*\mathcal O_{\mathbb P^1}(-1)$ 在 $C_i$ 上的限制为 $\mathcal O_{C_i}(-1)$，$i=1,2$。正规化映射诱导短正合列
\[
0\longrightarrow f_1^*\mathcal O_{\mathbb P^1}(-1)
\longrightarrow\mathcal O_{C_1}(-1)\oplus\mathcal O_{C_2}(-1)
\longrightarrow f_1^*\mathcal O_{q_1}(-1)\longrightarrow0.
\]
取上同调，得到一维空间
\[
H^1(C_1\cup C_2,f_1^*\mathcal O_{\mathbb P^1}(-1))
\simeq H^0(\mathbb P^1,\mathcal O_{q_1}(-1)).
\]
$T$ 在 $\mathcal O_{\mathbb P^1}(-1)$ 于 $q_1$ 处的作用为 $-\rho_1$，所以权为 $-\lambda_1$。因此，障碍丛在 $f_1$ 处的两个权为 $-\lambda_1,-\lambda_1$。

与 $\Gamma_2$ 对应的固定点 $Z_{\Gamma_2}$ 处的计算完全类似。只需交换 $\lambda_0,\lambda_1$，得到
\begin{align*}
\text{切空间的权：}&\quad\lambda_0-\lambda_1,\;2(\lambda_0-\lambda_1),\\
\text{障碍丛的权：}&\quad-\lambda_0,\;-\lambda_0.
\end{align*}
\SourcePage{314}
最后考虑 $\Gamma_3$。其中有 2 个旗、2 个顶点和 1 条边。由定理~\ref{thm:fixed-graph-normal-euler}，
\begin{align*}
e_{\Gamma_3}^{\mathrm F}&=\frac1{(\lambda_0-\lambda_1)(\lambda_1-\lambda_0)},\\
e_{\Gamma_3}^{\mathrm v}&=\frac{(\lambda_1-\lambda_0)(\lambda_0-\lambda_1)}
{((\lambda_0-\lambda_1)/2)((\lambda_1-\lambda_0)/2)},\\
e_{\Gamma_3}^{\mathrm e}&=\frac{2^2(\lambda_0-\lambda_1)^4}{2^4}.
\end{align*}
所以 $\operatorname{Euler}_T(N_{\Gamma_3})=(\lambda_0-\lambda_1)(\lambda_1-\lambda_0)$。也可像 $\Gamma_1$ 那样，通过显式计算得到这一结果。

还需要固定点 $\Gamma_3$ 处障碍丛的权。这里的固定点由二次映射 $f_3:\mathbb P^1\to\mathbb P^1$，$f_3(z_0,z_1)=(z_0^2,z_1^2)$ 表示。沿用 $\Gamma_1$ 的方法，只须计算 $T$ 在 $H^1(\mathbb P^1,f_3^*\mathcal O_{\mathbb P^1}(-1))$ 上作用的权，再将结果重复一次。

考虑由 Čech 上闭链
\[
\frac1{z_0z_1}\in Z^1(U_0\cap U_1,f_3^*\mathcal O_{\mathbb P^1}(-1))
\]
表示的非零上同调类，其中 $U_i=\{z\in\mathbb P^1:z_i\ne0\}$。变量 $x_i\in H^0(\mathbb P^1,\mathcal O_{\mathbb P^1}(1))$ 的权为 $\lambda_i$；由 $f_3$ 有 $x_i=z_i^2$，所以 $z_i$ 的权为 $\lambda_i/2$。于是 $1/(z_0z_1)$ 的权为 $-(\lambda_0+\lambda_1)/2$。因此，障碍丛的两个权均为 $-(\lambda_0+\lambda_1)/2$。

施加局部化公式所需的所有量现已齐备。积分
\[
\int_{\overline M_{0,0}(\mathbb P^1,2)}\operatorname{Euler}
\bigl(R^1\pi_{1*}e_1^*(\mathcal O_{\mathbb P^1}(-1)\oplus\mathcal O_{\mathbb P^1}(-1))\bigr)
\]
等于
\begin{align*}
&\frac{\lambda_1\lambda_1}{2(\lambda_1-\lambda_0)(2(\lambda_1-\lambda_0))}
+\frac{\lambda_0\lambda_0}{2(\lambda_0-\lambda_1)(2(\lambda_0-\lambda_1))}\\
&\qquad+\frac{(- (\lambda_0+\lambda_1)/2)(-(\lambda_0+\lambda_1)/2)}
{2(\lambda_1-\lambda_0)(\lambda_0-\lambda_1)}=\frac18.
\end{align*}
注意，每项分母中额外的因子 $2$，来自每个固定点处的阶为 $2$ 的自同构群，见例题~\ref{ex:fixed-graphs-p1-degree2}。这就完成了 $d=2$ 情形的证明。

对于一般的 $d$，\SourceCite{Manin2} 给出了一种证明，将上述方法与巧妙的组合求和结合起来。这里采用 R. Pandharipande 的简化方法证明任意 $d$ 的情形。其思路，是改变 $\mathcal O_{\mathbb P^1}(-1)\oplus\mathcal O_{\mathbb P^1}(-1)$ 上的环面作用，以简化计算。先将环面限制到子环面 $T'=\mathbb C^*\times\{1\}\subset T=(\mathbb C^*)^2$，记 $\lambda_0$ 在 $T'$ 上的限制为 $h$。在 $\overline M_{0,0}(\mathbb P^1,d)$ 上使用限制所得的作用，但在两个负线丛的直和上使用另一个 $T'$ 作用。

具体地，赋予 $\mathcal O_{\mathbb P^1}(-1)\oplus\mathcal O_{\mathbb P^1}(-1)$ 如下的作用：
\[
t\cdot(l(x_0,x_1),m(x_0,x_1))=(l(tx_0,x_1),m(x_0,t^{-1}x_1)).
\]
这里 $l,m$ 关于 $(x_0,x_1)$ 都是 $-1$ 次齐次的。注意，这个作用与 $T'$ 在 $\mathbb P^1$ 上的作用相容。
\SourcePage{315}
现在利用该作用和推论~\ref{cor:orbifold-localization-integration}，计算障碍丛的 Euler 类在 $\overline M_{0,0}(\mathbb P^1,d)$ 上的积分。容易看出，$T'$ 与 $T$ 作用的固定点集相同。设图 $\Gamma$ 对应一个固定点分支 $\overline M_\Gamma$。我们断言：若 $\Gamma$ 不止一条边，则它在局部化公式中的贡献为零。

证明并不困难。设 $f:C\to\mathbb P^1$ 是该分支中的一般稳定映射，将 $C$ 分解为分支 $C_i$，结点集记为 $\{r_j\}$。若图不止一条边，则至少有一个结点。正规化正合列为
\begin{align*}
0\longrightarrow f^*(\mathcal O_{\mathbb P^1}(-1)\oplus\mathcal O_{\mathbb P^1}(-1))
&\longrightarrow\bigoplus_i(f|_{C_i})^*(\mathcal O_{\mathbb P^1}(-1)\oplus\mathcal O_{\mathbb P^1}(-1))\\
&\longrightarrow\bigoplus_j f^*(\mathcal O_{\mathbb P^1}(-1)\oplus\mathcal O_{\mathbb P^1}(-1))_{r_j}
\longrightarrow0.
\end{align*}
取上同调，结点处纤维截面的空间，模去各收缩分支上的常值截面的像后，嵌入 $H^1(C,f^*(\mathcal O_{\mathbb P^1}(-1)\oplus\mathcal O_{\mathbb P^1}(-1)))$。

由于 $f\in\overline M_\Gamma$，每个 $f(r_j)$ 都是 $q_0$ 或 $q_1$。若 $f(r_j)=q_1$，则 $1/x_1$ 是负线丛在 $q_1$ 处的局部基，所以 $T'$ 在第一个纤维因子上的作用平凡。类似地，若 $f(r_j)=q_0$，以 $1/x_0$ 为局部基可知，$T'$ 在第二个纤维因子上的作用平凡。具有至少两条边的树必有价数至少为 $2$ 的顶点，因而上述商中至少保留一个零权。这一零权乘入障碍丛的等变 Euler 类的限制，故固定点分支 $\overline M_\Gamma$ 的贡献为零，断言得证。

因此，只需考虑对应一个 $d$ 次映射的单边图：
\begin{figure}[H]
\centering\includegraphics[width=.103\textwidth]{multiple-cover-graph-edge.png}
\end{figure}
对于该图，$\overline M_\Gamma$ 只有一点，自同构群为 $A_\Gamma=\mathbb Z/d\mathbb Z$；相应映射仍为 $f(z_0,z_1)=(z_0^d,z_1^d)$。

使用前面 $d=2$ 时的开覆盖 $U_0,U_1$，计算障碍丛的限制的权。一般的 $d$ 情形中，$H^1(\mathbb P^1,f^*(\mathcal O_{\mathbb P^1}(-1)\oplus\mathcal O_{\mathbb P^1}(-1)))$ 的一组基可由下列 Čech 上闭链给出：
\begin{equation}\label{eq:multiple-cover-cech-basis}
\left(\frac1{z_0^iz_1^{d-i}},0\right),\quad
\left(0,\frac1{z_0^iz_1^{d-i}}\right),\qquad1\le i\le d-1.
\end{equation}
$x_0,x_1$ 的权为 $h,0$；由 $f$ 有 $x_i=z_i^d$，故 $z_0,z_1$ 的权为 $h/d,0$。于是式~\eqref{eq:multiple-cover-cech-basis} 中两类上闭链的 $T'$ 权分别为 $-ih/d$ 和 $(d-i)h/d$，$1\le i\le d-1$。因此，障碍丛的等变 Euler 类在局部化公式中贡献
\[
(-1)^{d-1}((d-1)!)^2h^{2d-2}/d^{2d-2}.
\]
\SourcePage{316}
还须计算法丛 $N_\Gamma$ 的等变 Euler 类。在定理~\ref{thm:fixed-graph-normal-euler} 中取 $\lambda_0=h,\lambda_1=0$，得到
\begin{align*}
e_\Gamma^{\mathrm F}&=\frac1{(h-0)(0-h)}=-h^{-2},\\
e_\Gamma^{\mathrm v}&=\frac{(h-0)(0-h)}{((h-0)/d)((0-h)/d)}=d^2,\\
e_\Gamma^{\mathrm e}&=\frac{(-1)^d(d!)^2(h-0)^{2d}}{d^{2d}}
=\frac{(-1)^d(d!)^2h^{2d}}{d^{2d}}.
\end{align*}
由此可知
\[
\operatorname{Euler}_{T'}(N_\Gamma)=(-1)^{d-1}(d!)^2h^{2d-2}/d^{2d-2}.
\]
最后，因为处理的是代数叠，还要计入 $\overline M_\Gamma$ 的自同构因子 $a_\Gamma$。正合列~\eqref{eq:fixed-graph-automorphisms} 说明，这里 $a_\Gamma=d$，故分母中额外出现因子 $d$。综合以上结果，推论~\ref{cor:orbifold-localization-integration} 给出所求积分
\[
\frac{(-1)^{d-1}((d-1)!)^2h^{2d-2}/d^{2d-2}}
{d\cdot(-1)^{d-1}(d!)^2h^{2d-2}/d^{2d-2}}=\frac1{d^3}.
\]
这就证明了定理。
\end{proof}

\SourceCite[推论 3.5]{LLY} 给出了定理~\ref{thm:multiple-cover-localization} 的完全不同的证明。其思想与\hyperref[sec:quintic-mirror-theorem]{第~\ref*{sec:quintic-mirror-theorem} 节}所述镜像定理的证明相似，例题~\ref{ex:lly-multiple-cover-proof} 将介绍这一证明。

\subsection{六结点五次曲线的二重覆盖与 \texorpdfstring{$n_{10}$}{n10}}\label{subsec:nodal-quintic-double-covers}
\hyperref[subsec:instanton-numbers]{第~\ref*{subsec:instanton-numbers} 小节}研究五次三维簇 $V\subset\mathbb P^4$ 上的有理曲线时，已说明：当 $d\le9$ 时，瞬子数 $n_d$ 就是 $V$ 上次数为 $d$ 的有理曲线数。论证中的关键，是定理~\ref{thm:multiple-cover-localization}，它告诉我们刚性嵌入的光滑有理曲线的多重覆盖，如何对 $N_d=\langle I_{0,0,d}\rangle$ 作出贡献。不过，R. Pandharipande 向我们指出，该定理并不适用于 $V$ 上次数为 $5$ 的结点有理曲线的二重覆盖。这与 \SourceCite{Katz3,BL} 所发现的结点曲线多重覆盖的异常现象相似。下面将证明，每条这样的曲线对 $N_{10}=\langle I_{0,0,10}\rangle$ 的贡献为 $6+\frac18$，而非简单套用该定理所得的 $2^{-3}=\frac18$。\index{结点曲线!二重覆盖|hyperpage}

先解释为什么 $V$ 含有次数为 $5$ 的结点有理曲线。若 $P\subset\mathbb P^4$ 是平面，则 $C=V\cap P$ 是算术亏格为 $6$ 的平面五次曲线。我们希望在这些 $C$ 中找出有理曲线。若 $C$ 有 $6$ 个结点，它便是有理曲线。注意，$C$ 在 $p$ 处有普通结点，当且仅当 $P$ 在 $p$ 处与 $V$ 相切，且相切的二次项非退化。要求 $P$ 与 $V$ 相切，是关于 $P$ 的一个条件，故在 $6$ 个点处相切应给出 $6$ 个条件。曲线 $C$ 由射影平面 $P\subset\mathbb P^4$ 的 Grassmann 簇 $G(3,5)$ 参数化，而该簇的维数为 $6$，所以预期一般五次三维簇包含有限且非零数目的六结点有理平面五次曲线。这种朴素的维数计数可以严格化；已知其带重数的曲线数为 $17\,601\,000$ \SourceCite{Vainsencher}；这一计数的完全严格证明见 S.~Kleiman 与 R.~Piene, \emph{Node polynomials for families: methods and applications}, Math.\ Nachr.\ \textbf{271} (2004), 69--90。下面将证明，一般情形下所有重数都为 $1$，从而一般五次三维簇恰有 $17\,601\,000$ 条六结点平面五次曲线。

设 $C_n$ 是一般五次三维簇 $V$ 上这样的一条结点有理平面五次曲线。$C_n$ 的一个\emph{$d$ 次多重覆盖}，是指亏格零稳定映射 $f:C\to C_n$，满足 $f_*[C]=d[C_n]$，在 $V$ 中的次数为 $5d$。
\SourcePage{317}
也可将它看作通过 $C_n$ 分解的映射 $f:C\to V$。我们要找出这些多重覆盖对 Gromov--Witten 不变量 $N_{5d}$ 的贡献。对定理~\ref{thm:multiple-cover-localization} 的公式作朴素解释，会预期贡献为 $1/d^3$；我们将说明，这在 $d=2$ 时已不成立。论证分三步：
\begin{itemize}
\item 先计算 $C_n$ 的正规化映射的法丛。
\item 再证明 $C_n$ 的二重覆盖构成 $\overline M_{0,0}(V,10)$ 的若干连通分支的并。
\item 最后证明虚拟基本类 $[\overline M_{0,0}(V,10)]^{\mathrm{virt}}$ 限制到这些分支后的总次数为 $6+\frac18$。
\end{itemize}
先处理第一步，它将说明 $C_n$ 刚性嵌入 $V$。\hyperref[subsec:instanton-numbers]{第~\ref*{subsec:instanton-numbers} 小节}已使用过这个事实。

\begin{lemma}\label{lem:nodal-quintic-normal-bundle}\index{法丛!结点曲线的正规化|hyperpage}
若 $V$ 是一般五次三维簇，$g:\mathbb P^1\to C_n$ 是 $C_n\subset V$ 的正规化，则 $g$ 在 $V$ 中的法丛为 $N_g\simeq\mathcal O_{\mathbb P^1}(-1)\oplus\mathcal O_{\mathbb P^1}(-1)$。
\end{lemma}
\begin{proof}
$N_g$ 是 $\mathbb P^1$ 上的秩二向量丛，所以存在 $a,b$，使 $N_g\simeq\mathcal O_{\mathbb P^1}(a)\oplus\mathcal O_{\mathbb P^1}(b)$。由于 $V$ 是 Calabi--Yau 簇，伴随公式给出 $a+b=-2$。因此，$a=b=-1$ 当且仅当 $h^0(\mathbb P^1,N_g)=0$。

下面将 \SourceCite[第 152--153 页]{Katz2} 的论证适用于此情形。令 $\mathbb P=\mathbb P(H^0(\mathbb P^4,\mathcal O_{\mathbb P^4}(5)))$ 为五次超曲面的射影参数空间，$M$ 为 $\mathbb P^4$ 中亏格零不可约有理平面五次曲线的模空间。$M$ 纤维化于 $G(3,5)$ 上，纤维同构于 $\overline M_{0,0}(\mathbb P^2,5)$ 的开子集，故 $M$ 不可约且维数为 $20$；六结点曲线构成其中的开集。令
\[
I=\{(D,V)\in M\times\mathbb P:D\subset V\}
\]
为关联对应。由 Riemann--Roch 定理，$h^0(D,\mathcal O_D(5))=20$。每个 $D\in M$ 都是 $\mathbb P^4$ 中的完全交，故为射影正规。于是可算出投影 $\pi_1:I\to M$ 的纤维是 $105$ 维射影空间。所以 $I$ 不可约，维数为 $125$，与 $\mathbb P$ 的维数相同。

现在令
\[
I_1=\{(D,V)\in I:h^0(\mathbb P^1,N_g)\ge1\},
\]
其中 $g:\mathbb P^1\to D\subset V$ 是正规化映射，在上述六结点开集及光滑 $V$ 上考察此法丛。由前面对 $N_g$ 的描述，$(D,V)\in I\setminus I_1$ 当且仅当 $N_g\simeq\mathcal O_{\mathbb P^1}(-1)\oplus\mathcal O_{\mathbb P^1}(-1)$。稍后将证明这个补集非空。假设如此，由 $I$ 的不可约性，便有 $\dim I_1\le124$，所以 $\pi_2(I_1)$ 的闭包是 $\mathbb P$ 的真子簇。对补集中的 $V$，引理的断言对 $V$ 上所有这类平面有理五次曲线都成立。

因此，只需给出一个例子。设 $(z_0,z_1)$ 为 $\mathbb P^1$ 的齐次坐标，考虑下列映射的像曲线 $D$：
\begin{equation}\label{eq:nodal-quintic-parametrization}
\begin{aligned}
(z_0,z_1)\longmapsto\bigl(&z_0^5+z_1^5,\;
 z_0z_1(z_0^3+z_0^2z_1+z_0z_1^2),\\
 &z_0z_1(z_0^3+3z_0^2z_1+2z_0z_1^2+z_1^3),\;0,\;0\bigr).
\end{aligned}
\end{equation}
曲线 $D$ 位于五次超曲面 $F=0$ 中，其中
\begin{equation}\label{eq:nodal-quintic-equation}
\begin{aligned}
F(x_0,\ldots,x_4)={}&x_0^3x_1^2+3x_0^2x_1^3+20x_0x_1^4+61x_1^5
 -x_0^3x_1x_2+x_0^2x_1^2x_2\\
&-24x_0x_1^3x_2-25x_1^4x_2-3x_0^2x_1x_2^2
 +2x_0x_1^2x_2^2-40x_1^3x_2^2\\
&+x_0^2x_2^3+4x_0x_1x_2^3+35x_1^2x_2^3
 -x_0x_2^4-10x_1x_2^4+x_2^5,
\end{aligned}
\end{equation}
\SourcePage{318}
$x_0,\ldots,x_4$ 是 $\mathbb P^4$ 的坐标。包含 $D$ 的一般五次超曲面的方程形如
\begin{equation}\label{eq:quintic-containing-nodal-curve}
F+x_3G+x_4H=0,
\end{equation}
其中 $G,H$ 都是四次齐次多项式。

由法丛正合列
\begin{equation}\label{eq:nodal-normalization-tangent-sequence}
0\longrightarrow T_{\mathbb P^1}\longrightarrow g^*T_V\longrightarrow N_g\longrightarrow0,
\end{equation}
可知，$h^0(\mathbb P^1,N_g)=0$ 当且仅当 $h^0(\mathbb P^1,g^*T_V)=3$。由 Riemann--Roch 定理，这又等价于 $H^1(\mathbb P^1,g^*T_V)=0$。我们略加推广 \SourceCite{Katz1} 中的计算方法来求这个上同调群。设 $f:C\to\mathbb P^4$，$C$ 为可能可约的亏格零有理曲线，令 $\mathcal L=f^*\mathcal O_{\mathbb P^4}(1)$。若 $\widetilde F=0$ 是 $V$ 的方程，记 $\widetilde F_j=f^*(\partial\widetilde F/\partial x_j)\in H^0(C,4\mathcal L)$。考虑正合列
\[
0\longrightarrow\mathcal O_{\mathbb P^4}\longrightarrow\mathcal O_{\mathbb P^4}(1)^5\longrightarrow T_{\mathbb P^4}\longrightarrow0
\]
及
\[
0\longrightarrow T_V\longrightarrow T_{\mathbb P^4}|_V\longrightarrow\mathcal O_V(5)\longrightarrow0.
\]
由 $f$ 拉回并结合这两个正合列，得到映射
\[
\Phi_V:H^0(C,\mathcal L)^5\longrightarrow H^0(C,5\mathcal L),\qquad
(s_0,\ldots,s_4)\longmapsto\sum_{j=0}^4s_j\widetilde F_j.
\]
简单的图表追踪给出 $H^1(C,f^*T_V)\simeq\operatorname{coker}\Phi_V$。这里用了 $h^1(C,\mathcal O_C)=h^1(C,f^*T_{\mathbb P^4})=0$；后一个等式来自 $\mathbb P^4$ 的凸性。

这样就能用线性代数和计算机直接求法丛。例如，取 $G=x_0^4,H=x_1^4$，$\Phi_V$ 的矩阵有整数元素；它是 $30\times26$ 矩阵，计算所得秩为 $26$。相应的五次超曲面有奇点，但形如式~\eqref{eq:quintic-containing-nodal-curve} 的一般五次超曲面光滑。由秩的半连续性，也能找到满足条件的光滑 $V$。
\end{proof}

该引理说明，一般 $V$ 上的每条六结点有理平面五次曲线的重数均为 $1$，所以 $V$ 恰有 $17\,601\,000$ 条这样的曲线。

下面处理二重覆盖研究的第二步，证明这些覆盖给出 $\overline M_{0,0}(V,10)$ 的 $7$ 个连通分支。令 $M_{2,C_n}(V)\subset\overline M_{0,0}(V,10)$ 为分类所有零标记亏格零稳定二重覆盖 $f:C\to C_n$ 的子空间。
\begin{proposition}\label{prop:nodal-double-cover-components}\index{稳定映射!二重覆盖|hyperpage}
$M_{2,C_n}(V)$ 由模空间 $\overline M_{0,0}(V,10)$ 的 $7$ 个连通分支构成，且：
\begin{enumerate}[label=(\roman*)]
\item 一个分支同构于 $\overline M_{0,0}(\mathbb P^1,2)$，由通过正规化映射分解的二重覆盖组成。
\item 其余 $6$ 个分支是 $\overline M_{0,0}(V,10)$ 的光滑点，对应二重覆盖 $f:C_1\cup C_2\to C_n$，其中 $f|_{C_i}$ 是正规化映射，且 $f$ 在 $C_n$ 的一个结点上方为局部同构。
\end{enumerate}
\end{proposition}
\SourcePage{319}
\begin{proof}
先考虑通过正规化映射 $g:\mathbb P^1\to C_n$ 分解的二重覆盖。令 $\mathcal M_1$ 为这类稳定二重覆盖的代数叠，并考虑相应的模空间 $M_1\subset\overline M_{0,0}(V,10)$。它们恰是可写成 $f=g\circ h$ 的映射，其中 $h:C\to\mathbb P^1$ 属于 $\overline M_{0,0}(\mathbb P^1,2)$。我们断言，$M_1$ 是一个分支。为此，在 $f\in M_1$ 处对 $\overline M_{0,0}(V,10)$ 作无穷小分析。由于 $\overline M_{0,0}(\mathbb P^1,2)$ 在轨形意义下光滑、完备且二维，只需证明环境模空间在 $f$ 处的切空间也为二维。照例，这个切空间是 $\operatorname{Ext}^1(f^*\Omega_V^1\to\Omega_C^1,\mathcal O_C)$。

先设 $(f,C)\in M_1$ 且 $C$ 不可约。那么 $C\simeq\mathbb P^1$，相应的 $\operatorname{Ext}^*$ 长正合列化为
\[
0\longrightarrow H^0(\mathbb P^1,T_{\mathbb P^1})\longrightarrow H^0(\mathbb P^1,f^*T_V)
\longrightarrow\operatorname{Ext}^1(f^*\Omega_V^1\to\Omega_{\mathbb P^1}^1,\mathcal O_{\mathbb P^1})\longrightarrow0,
\]
并给出同构 $\operatorname{Ext}^2(f^*\Omega_V^1\to\Omega_{\mathbb P^1}^1,\mathcal O_{\mathbb P^1})\simeq H^1(\mathbb P^1,f^*T_V)$。由引理~\ref{lem:nodal-quintic-normal-bundle} 和式~\eqref{eq:nodal-normalization-tangent-sequence}，
\[
g^*T_V\simeq\mathcal O_{\mathbb P^1}(2)\oplus\mathcal O_{\mathbb P^1}(-1)\oplus\mathcal O_{\mathbb P^1}(-1),
\]
从而
\[
f^*T_V\simeq\mathcal O_{\mathbb P^1}(4)\oplus\mathcal O_{\mathbb P^1}(-2)\oplus\mathcal O_{\mathbb P^1}(-2).
\]
立即可知上述 $\operatorname{Ext}^1$ 的维数为 $2$。

也可能有 $(f,C)\in M_1$，其中 $C$ 有两个分支，每个都以次数 $1$ 映满 $C_n$。此时同样可算出 $\operatorname{Ext}^1(f^*\Omega_V^1\to\Omega_C^1,\mathcal O_C)$ 为二维，细节略去。命题的第 (i) 项得证。

现在设 $f:C\to C_n$ 是不能通过正规化 $g$ 分解的二次覆盖。写 $C=\bigcup_iC_i$，记 $f$ 在 $C_i$ 上的次数为 $d_i$，则 $\sum_i d_i=2$。最多只有两个分支以正次数映到 $C_n$。若有一棵分支组成的树被 $f$ 收缩，即树上所有 $d_i=0$，它最多与另两个分支相遇，与稳定性矛盾。因此，$C$ 只能有一或两个分支。若只有一个，则 $C$ 正规，$f$ 必通过 $g$ 分解。所以 $C=C_1\cup C_2$ 是两个光滑有理分支的并，且 $d_1=d_2=1$。每个 $f|_{C_i}$ 因而可与正规化映射 $g$ 等同，只需描述如何将二者粘合。

若 $f$ 将 $C$ 的结点映到 $C_n$ 的光滑点，显然 $f$ 通过 $g$ 分解。若结点映到 $C_n$ 的结点，则 $f$ 通过 $g$ 分解，当且仅当 $C$ 在结点附近的两个分支映到 $C_n$ 在像结点处的同一分支。因此，不通过 $g$ 分解时，两个分支必须分别映到像结点的两个不同分支。换言之，$f$ 在该结点上方为局部同构。

于是，恰有 $6$ 个不通过正规化分解的二重覆盖，每个结点对应一个。记结点为 $p_1,\ldots,p_6$，记 $f_i:C=C_1\cup C_2\to C_n\subset V$ 为如下覆盖：在每个 $C_j$ 上为正规化映射，将 $C$ 的结点映到 $p_i$，并使 $C_1,C_2$ 分别映到 $p_i$ 处的两个不同分支。

以上说明，$M_{2,C_n}(V)$ 由这六个映射及通过正规化分解的二重覆盖分支 $\overline M_{0,0}(\mathbb P^1,2)$ 组成。为完成证明，只需说明 $f_i$ 是环境模空间的光滑点，等价地，其切空间 $\operatorname{Ext}^1(f_i^*\Omega_V^1\to\Omega_C^1,\mathcal O_C)$ 的维数为零。
\SourcePage{320}
由通常的正合列，只须证明映射
\begin{equation}\label{eq:nodal-cover-smoothing-surjectivity}
\operatorname{Ext}^1(\Omega_C^1,\mathcal O_C)\longrightarrow H^1(C,f_i^*T_V)
\end{equation}
满射。令 $K$ 为 $f_i^*\Omega_V^1\to\Omega_C^1$ 的核。用局部坐标容易看出，$K$ 是秩为 $2$ 的局部自由层；式~\eqref{eq:nodal-cover-smoothing-surjectivity} 的满射性等价于 $H^1(C,K^*)=0$。下面计算这个上同调群。

采用间接论证。考虑关联对应
\begin{equation}\label{eq:nodal-quintic-point-incidence}
\mathcal I=\{(p,D,V)\in\mathbb P^4\times M\times\mathbb P:
 p\text{ 是 }D\subset V\text{ 的一个结点}\},
\end{equation}
其中 $M$ 与引理~\ref{lem:nodal-quintic-normal-bundle} 的证明中一样，是亏格零不可约有理平面五次曲线的模空间。

在固定点处有结点的不可约有理曲线组成的集合不可约。这是因为，例如在 $(1,0,0,0,0)$ 处有结点的曲线，可以用如下映射参数化：
\[
\begin{aligned}
(z_0,z_1)\longmapsto\bigl(&\alpha_0(z_0,z_1),\;z_0z_1\alpha_1(z_0,z_1),\;z_0z_1\alpha_2(z_0,z_1),\\
 &z_0z_1\alpha_3(z_0,z_1),\;z_0z_1\alpha_4(z_0,z_1)\bigr),
\end{aligned}
\]
这里 $\alpha_0$ 为五次齐次式，$\alpha_i$ 为三次齐次式，$1\le i\le4$。在这参数化中，任意将 $(1,0)$、$(0,1)$ 选为映到结点的两点。式~\eqref{eq:nodal-quintic-parametrization} 是其特例。所有 $\alpha_i$ 的参数空间不可约；限制到平面曲线时，再选取包含该点的平面及其中的三个坐标。平面的参数空间亦不可约，故所需结论成立。

利用这一点，可证明式~\eqref{eq:nodal-quintic-point-incidence} 定义的 $\mathcal I$ 不可约：先投影到 $\mathbb P^4\times M$，再投影到 $\mathbb P^4$，应用上一段的结论。

此后，引理~\ref{lem:nodal-quintic-normal-bundle} 的证明所用的论证，将问题归结为找出一个 $(p,D,V)\in\mathcal I$，满足 $H^1(C,K^*)=0$。取 $D$ 为式~\eqref{eq:nodal-quintic-equation} 给出的曲线，$p=(1,0,0,0,0)$。由参数化~\eqref{eq:nodal-quintic-parametrization}，该结点是 $(1,0)$ 和 $(0,1)$ 的共同像。用 Macaulay \SourceCite{BS}，可将 $C$ 显式构造为 $\mathbb P^2$ 中的两条直线，将 $f$ 构造为显式映射 $\mathbb P^2\to\mathbb P^4$ 在 $C$ 上的限制，再追踪上述构造，得到层 $K^*$。计算给出所需的 $H^1(C,K^*)=0$。虽然 Macaulay 在非零特征下计算，但由上同调的上半连续性，仍可推出一般情形下的消失。
\end{proof}

最后，计算虚拟基本类 $[\overline M_{0,0}(V,10)]^{\mathrm{virt}}$ 限制到命题~\ref{prop:nodal-double-cover-components} 的各分支后的次数。
\begin{theorem}\label{thm:nodal-quintic-double-cover-contribution}\index{五次三维簇!六结点五次曲线|hyperpage}
若 $C_n$ 是一般五次三维簇 $V\subset\mathbb P^4$ 上的六结点有理平面五次曲线，则 $C_n$ 的二重覆盖对 Gromov--Witten 不变量 $N_{10}=\langle I_{0,0,10}\rangle$ 的贡献为 $6+\frac18$。
\end{theorem}
\begin{proof}
先考虑通过正规化分解的二重覆盖所构成的分支 $\overline M_{0,0}(\mathbb P^1,2)$。直接沿用定理~\ref{thm:multiple-cover-localization} 的证明即可；切空间与障碍空间的权均不变，故该分支对 $N_{10}$ 的贡献为 $1/8$。

再看由 $f_1,\ldots,f_6$ 表示的六个分支。关键在于，它们都是具有预期维数的光滑点。由定义~\ref{def:convex-variety} 后的讨论，虚拟基本类在各点的限制为通常的点的基本类，即系数为 $1$ 的点。因此，六个分支各贡献 $1$，定理得证。
\end{proof}

\SourcePage{321}
\section{等变 Gromov--Witten 不变量}\label{sec:equivariant-gw}\index{Gromov--Witten不变量@Gromov--Witten 不变量!等变不变量|hyperpage}
现在简要讨论等变 Gromov--Witten 不变量。为简便起见，先考虑 $g=0$ 且 $X$ 为凸簇的情形。

\subsection{凸簇的亏格零不变量}\label{subsec:convex-equivariant-gw}
设 $X$ 光滑且凸。如\hyperref[subsec:computing-virtual-class]{第~\ref*{subsec:computing-virtual-class} 小节}所述，这意味着 $[\overline M_{0,n}(X,\beta)]^{\mathrm{virt}}=[\overline M_{0,n}(X,\beta)]$。于是，可以如下定义亏格零等变 Gromov--Witten 不变量。

设代数群 $G$ 在 $X$ 上作用，并保持曲线类 $\beta$。则 $h\in G$ 按公式 $h\cdot f=\tau_h\circ f$ 自然地作用于 $\overline M_{0,n}(X,\beta)$，其中 $\tau_h:X\to X$ 是与 $h$ 对应的自同构。在\hyperref[sec:stable-map-localization]{第~\ref*{sec:stable-map-localization} 节}中，已经对 $X=\mathbb P^r$、$G=(\mathbb C^*)^{r+1}$ 的情形详细分析了这个作用。

由于 $[\overline M_{0,n}(X,\beta)]^{\mathrm{virt}}=[\overline M_{0,n}(X,\beta)]$，\hyperref[subsec:gw-algebraic-classes]{第~\ref*{subsec:gw-algebraic-classes} 小节}给出的 Gromov--Witten 不变量定义简化为
\[
\langle I_{0,n,\beta}\rangle(\alpha_1,\ldots,\alpha_n)
=\int_{\overline M_{0,n}(X,\beta)}e_1^*(\alpha_1)\cup\cdots\cup e_n^*(\alpha_n),
\]
其中 $\alpha_i\in H^*(X)$，而
\[
e_i:\overline M_{0,n}(X,\beta)\longrightarrow X
\]
是\IndexLocation{indexloc:0276:321}赋值映射，将稳定映射 $f:(C,p_1,\ldots,p_n)\to X$ 映为 $f(p_i)$（见\hyperref[subsec:moduli-to-gw]{第~\ref*{subsec:moduli-to-gw} 小节}）。

要给出它的等变版本，先注意到，对于 $G$ 在 $\overline M_{0,n}(X,\beta)$ 和 $X$ 上的作用，$e_i$ 是等变的。因此有等变拉回映射
\[
e_i^*:H_G^*(X)\longrightarrow H_G^*(\overline M_{0,n}(X,\beta)).
\]
此外，$X$ 为凸簇，所以 $\overline M_{0,n}(X,\beta)$ 是轨形，这意味着存在等变积分
\[
\int_{\overline M_{0,n}(X,\beta)_G}:
H_G^*(\overline M_{0,n}(X,\beta))\longrightarrow H^*(BG).
\]
定义所需的各个要素现在都已具备。
\begin{definition}\label{def:equivariant-gw}
设 $X$ 光滑且凸，$\alpha_1,\ldots,\alpha_n\in H_G^*(X)$。则表达式
\[
\langle I_{0,n,\beta}\rangle_G(\alpha_1,\ldots,\alpha_n)
=\int_{\overline M_{0,n}(X,\beta)_G}e_1^*(\alpha_1)\cup\cdots\cup e_n^*(\alpha_n)
\in H^*(BG)
\]
称为\emph{亏格零等变 Gromov--Witten 不变量}。
\end{definition}

这个定义与通常的 Gromov--Witten 不变量有如下联系。由\hyperref[subsec:equivariant-cohomology]{第~\ref*{subsec:equivariant-cohomology} 小节}，有“\IndexLocation{indexloc:0538:321}非等变极限”映射 $i_{\mathrm{point}}^*:H^*(BG)\to H^*(\mathrm{point})=\mathbb C$，以及 $i_X^*:H_G^*(X)\to H^*(X)$。于是，恒等式
\[
i_{\mathrm{point}}^*\langle I_{0,n,\beta}\rangle_G(\alpha_1,\ldots,\alpha_n)
=\langle I_{0,n,\beta}\rangle(i_X^*\alpha_1,\ldots,i_X^*\alpha_n)
\]
可由\hyperref[subsec:chern-polynomial-localization]{第~\ref*{subsec:chern-polynomial-localization} 小节}的交换图~\eqref{eq:equivariant-integral-forgetting-diagram} 证明。因此，等变 Gromov--Witten 不变量是通常 Gromov--Witten 不变量的精细化，包含来自群作用的更多信息。

接下来，设 $E$ 是光滑凸簇 $X$ 上的\IndexLocation{indexloc:0255:321}等变向量丛。如果光滑子簇 $Y\subset X$ 由 $E$ 的一个一般
\SourcePage{322}
截面的零点集定义，那么即使 $Y$ 本身没有群作用，也能为 $Y$ 的某些 Gromov--Witten 不变量定义等变版本。这里不讨论一般情形，只说明一个对我们尤其重要的例子。
\SourceExampleNumber{9.3.1}{1}
\begin{example}\label{ex:equivariant-quintic-gw}
设 $V$ 为五次三维簇。由例题~\ref{ex:quintic-gw-euler-integral}，有 Gromov--Witten 不变量
\[
N_d=\langle I_{0,0,d}\rangle
=\int_{\overline M_{0,0}(\mathbb P^4,d)}\operatorname{Euler}(\mathcal V_d),
\]
其中 $\mathcal V_d$ 是 $\overline M_{0,0}(\mathbb P^4,d)$ 上的向量丛，其在 $(f,C)\in\overline M_{0,0}(\mathbb P^4,d)$ 处的纤维为 $H^0(C,f^*\mathcal O_{\mathbb P^4}(5))$。由于 $T=(\mathbb C^*)^5$ 自然地作用于 $\overline M_{0,0}(\mathbb P^4,d)$ 和 $\mathcal V_d$，这个不变量的等变版本为
\[
\langle I_{0,0,d}\rangle_T
=\int_{\overline M_{0,0}(\mathbb P^4,d)_T}\operatorname{Euler}_T(\mathcal V_d)
\in H^*(BT).
\]
对于例题~\ref{ex:quintic-quartics-localization} 中 $N_4$ 的计算，可以这样理解：先应用推论~\ref{cor:orbifold-localization-integration} 计算等变 Gromov--Witten 不变量 $\langle I_{0,0,4}\rangle_T$，再取非等变极限，得到所需的 $N_4$。
\end{example}

以下不会显式使用等变 Gromov--Witten 不变量，尽管\hyperref[chap:mirror-theorem]{第十一章}的一些公式隐含着它们。不过，\hyperref[chap:mirror-theorem]{第十一章}确实会采用上述策略：先处理某些等变类，再在非等变极限中得到有意义的结果。

\subsection{虚拟基本类的局部化}\label{subsec:virtual-localization}\index{虚拟基本类!局部化|hyperpage}
要为任意光滑射影簇 $X$ 和任意亏格 $g\ge0$ 定义等变 Gromov--Witten 不变量，需要虚拟基本类的等变版本。当 $X$ 具有 $\mathbb C^*$ 作用时，\SourceCite{GPa} 已利用 $X$ 的等变 Chow 群 $A_*^{\mathbb C^*}(X)$ 建立了这一理论。

这个理论有两个关键点。第一个，正如所预期的，是定义等变虚拟基本类
\[
[\overline M_{g,n}(X,\beta)]_{\mathbb C^*}^{\mathrm{virt}}
\in A_*^{\mathbb C^*}(\overline M_{g,n}(X,\beta)).
\]
由此容易得到等变 Gromov--Witten 类。第二个关键要素是局部化公式。这也不意外，因为已经看到，局部化在计算 Gromov--Witten 不变量的贡献中发挥着重要作用。

基本思想是，\SourceCite{BF,Behrend} 的构造将等变虚拟基本类
\[
[Y]_{\mathbb C^*}^{\mathrm{virt}}\in A_*^{\mathbb C^*}(Y)
\]
赋予任何带有 $\mathbb C^*$ 作用和 $\mathbb C^*$ 等变完美障碍理论的代数概形 $Y$；相应的构造也适用于 Deligne--Mumford 叠。将它应用于 $Y=\overline M_{g,n}(X,\beta)$ 时，由 $X$ 上的作用诱导 $\mathbb C^*$ 作用，而定义通常虚拟基本类所用的完美障碍理论也自然继承 $\mathbb C^*$ 作用，从而给出
\[
[\overline M_{g,n}(X,\beta)]_{\mathbb C^*}^{\mathrm{virt}}
\in A_*^{\mathbb C^*}(\overline M_{g,n}(X,\beta)).
\]

此外，\SourceCite{GPa} 说明，等变虚拟基本类有很好的局部化性质。在上述 $Y$ 带有 $\mathbb C^*$ 作用及 $\mathbb C^*$ 等变完美障碍理论的情形，固定点轨迹的分支 $Z_j$ 继承作用权为零的完美障碍理论，因而得到 $A_*^{\mathbb C^*}(Z_j)$
\SourcePage{323}
中的等变虚拟基本类。$Z_j$ 的\emph{虚法丛} $N_j^{\mathrm{virt}}$ 来自障碍理论所确定的虚切空间的“移动部分”。于是，等变虚拟基本类的局部化公式 \SourceCite{GPa} 为
\[
[Y]_{\mathbb C^*}^{\mathrm{virt}}
=\sum_j i_{j!}\left(\frac{[Z_j]_{\mathbb C^*}^{\mathrm{virt}}}
{\operatorname{Euler}_{\mathbb C^*}(N_j^{\mathrm{virt}})}\right),
\]
其中，如往常一样，$\operatorname{Euler}_{\mathbb C^*}(N_j^{\mathrm{virt}})$ 是 $N_j^{\mathrm{virt}}$ 的等变 Euler 类。

将这一公式推广到 $G=(\mathbb C^*)^n$ 是直接的。这个局部化公式在定理~\ref{thm:givental-toric-mirror} 中起着关键作用。\SourceCite{GPa} 所述的局部化方法对 Gromov--Witten 不变量还有一些很好的应用，包括例题~\ref{ex:quintic-genus-one} 中提到的亏格一 Gromov--Witten 不变量的多重覆盖贡献。
